% !TeX root = ../../Book5.tex
% 纲要标题：## 第8章　消解、Koszul 复形与深度
% 纲要位置：几何学计划纲要.md，第 5299 行。
% 正文状态：初稿；尚未正式审核。

% 纲要内容开始
% 
% > 本章目标与证明范围：局部同调工具自成一章，实际计算非横截交和深度。先建立 Ext、Tor 与维数平移，再使用 Koszul 正合性和深度引理，不把抽象函子当作已学黑箱。
% 
% 纲要内容结束

\chapter{消解、Koszul 复形与深度}
\label{g5:ch:08}
\BookChapterNumber{g5:ch:08}

\begin{chapterabstract}
	本章从自由消解出发，构造 Ext、Tor 及其长正合列，用同调记录张量和 Hom 不能保持正合的位置。Koszul 复形把一组方程组织成有限复形；正则序列给出它的正合性，局部有限性则允许从同调消失反推正则性。随后证明正则序列长度、Ext 首次非零次数与深度之间的等价，建立深度引理和 Auslander–Buchsbaum 公式。最后计算平面曲线的交、完全交与嵌入点，比较支集相同而局部同调不同的模型。
\end{chapterabstract}

\begin{learninggoals}
	\begin{itemize}
		\item 逐层构造自由消解，证明链提升存在且在链同伦意义下唯一，并据此定义 Ext、Tor。
		\item 构造长正合列中的连接映射，用维数平移及显式消解计算导出障碍。
		\item 写出 Koszul 微分，证明正则序列与局部有限模上 Koszul 正合性的等价，检查末端非零条件及基变换限制。
		\item 用关联素理想和 Ext 证明深度的不同刻画、深度与维数的比较以及深度引理，准确表述模的 \(S_r\) 条件。
		\item 构造极小自由消解，从剩余域张量读取 Betti 数与投射维数，并完整证明 Auslander–Buchsbaum 公式。
		\item 计算横截、相切、自交及嵌入点模型，说明交的长度、Tor、深度分别检测什么。
	\end{itemize}
\end{learninggoals}

在 \(R=k[x,y]_{(x,y)}\) 中，乘 \(x\) 给出短正合列
\[
	0\longrightarrow R\xrightarrow{x}R\longrightarrow R/(x)\longrightarrow0.
\]
若与 \(R/(y)\) 张量，左端的乘 \(x\) 仍单射；若与 \(R/(x)\) 张量，同一张映射却变为零。坐标轴相交和直线自交因而产生了不同的正合性行为。第六章用平坦性判断一个模是否对所有短正合列都保持单射；现在要对给定的两个模保留失败的部分，并继续追踪更高阶的关系。

\ProseParagraphBreak

自由消解将模写成自由模之间连续的生成与关系。同伦比较保证换一份消解不会改变所得同调。对一组方程，Koszul 复形提供统一而有限的候选消解；这些方程逐次成为非零因子时，候选消解确实正合。局部环中能持续选出多少个这样的方程，就是深度所衡量的内容。极小消解又从另一端测量需要多少层关系，有限投射维数时两种计数由一个等式连接。

\ProseParagraphBreak

本章使用第二章的 Noether 模、关联素理想、投射模与 Nakayama 引理，第三章的素理想避开和维数理论，第六章的平坦判据；完备化只用于章末的一道习题。主要参考为\cite[附录B，第274—280、282页；第16节；第19节，第153—155页]{Matsumura1989}。各项同调工具均在本章建立，后续局部环的结构判据与交理论将调用这里的结论。

% !TeX root = ../../Book5.tex
% 纲要标题：### 8.1 自由消解与导出障碍
% 纲要位置：几何学计划纲要.md，第 5303 行。
% 正文状态：初稿；尚未正式审核。

% 纲要内容开始
% 
% * 自由消解的构造
% * 比较定理和链同伦
% * Ext、Tor 与长正合列
% * 维数平移和合冲模
% * 非平坦张量的显式 Tor
% 
% 纲要内容结束

\section{自由消解与导出障碍}
\label{g5:sec:0801}

第二章已经看到，张量积可能把一个单射变成零映射，\(\operatorname{Hom}\) 也可能使满射失去满性。自由模上的计算却容易控制：映射由基向量的像决定，张量自由模则只是取直和。因而可以先用自由模表示一个模，再继续用自由模表示关系、关系之间的关系。这样得到的消解把原来隐藏在核中的信息分层记录；张量或取 \(\operatorname{Hom}\) 之后，留下的同调便是这些运算失去正合性的具体障碍。

\subsection{自由消解的构造}
\label{g5:subsec:0801:01}

一个呈示 \(F_1\to F_0\to M\to0\) 只说明 \(M\) 的生成元及其关系。第一张映射也可能有核，因此关系本身还可能满足关系。为了不断记录下去，先约定复形与同调的记号。

\begin{definition}\label{g5:def:module-complex}
	设 \(A\) 为交换含幺环。一个 \(A\) 模的\term[lianfuxing]{链复形}[chain complex]是模与同态组成的序列
	\[
		C_\bullet:\quad\cdots\longrightarrow C_{i+1}
		\xrightarrow{\mathrm{d}_{i+1}}C_i\xrightarrow{\mathrm{d}_i}C_{i-1}
		\longrightarrow\cdots,
		\qquad \mathrm{d}_i \mathrm{d}_{i+1}=0.
	\]
	称 \(\mathrm{d}_i\) 为微分，称
	\[
		Z_i(C)=\ker \mathrm{d}_i,\qquad B_i(C)=\operatorname{im}\mathrm{d}_{i+1}
	\]
	中的元素分别为\term[xunhuan]{循环}[cycle]与\term[bianjie]{边界}[boundary]。其第 \(i\) 个\term[tongdiaomo]{同调模}[homology module]为
	\[
		H_i(C_\bullet)=Z_i(C)/B_i(C).
	\]
	若箭头使次数增加，即 \(\mathrm{d}^i:C^i\to C^{i+1}\)，则称为\term[shanglianfuxing]{上链复形}[cochain complex]，并定义其\term[shangtongdiaomo]{上同调模}[cohomology module]
	\[
		H^i(C^\bullet)=\ker \mathrm{d}^i/\operatorname{im}\mathrm{d}^{i-1}.
	\]
\end{definition}

条件 \(\mathrm{d}^2=0\) 保证边界是循环，故上述商有定义。\(H_i=0\) 正好表示复形在 \(C_i\) 处正合：所有被下一张映射消去的元素都已经来自上一项。下标与上标只是箭头方向的约定，令 \(C_i=C^{-i}\) 就能将两种记法互换。本节的消解都放在非负下标次数，其余次数的模取零。

\begin{definition}\label{g5:def:free-resolution}
	设 \(A\) 为交换含幺环，\(M\) 为 \(A\) 模。一个\term[ziyouxiaojie]{自由消解}[free resolution]是正合列
	\[
		\cdots\longrightarrow F_2\xrightarrow{\mathrm{d}_2}F_1
		\xrightarrow{\mathrm{d}_1}F_0\xrightarrow{\varepsilon}M
		\longrightarrow0,
	\]
	其中每个 \(F_i\) 都为自由 \(A\) 模。将不含 \(M\) 的复形记为 \(F_\bullet\)，将 \(\varepsilon\) 称为\term[zengguangyingshe]{增广映射}[augmentation map]。若将自由模换成投射模，则称为\term[toushexiaojie]{投射消解}[projective resolution]。
\end{definition}

于是 \(H_i(F_\bullet)=0\) 对 \(i>0\) 成立，而增广给出 \(H_0(F_\bullet)\cong M\)。消解中最右端的 \(M\) 与计算同调时的零次项 \(F_0\) 是不同对象；保留这个区分，可以避免次数错位。

\begin{proposition}\label{g5:prop:free-resolution-exists}
	任意交换含幺环 \(A\) 上的任意模 \(M\) 都有自由消解。若 \(A\) Noether 且 \(M\) 有限生成，则可以使消解中每个自由模都具有有限秩。
\end{proposition}
\begin{proof}
	选 \(M\) 的一组生成元；不要求有限，例如可取其全部元素。以这些生成元为基符号构造自由模 \(F_0\)，将基符号送到相应元素，得到满射 \(\varepsilon:F_0\to M\)。再选 \(\ker\varepsilon\) 的生成元，以它们构造自由模 \(F_1\) 及满射 \(F_1\to\ker\varepsilon\)，与包含复合得到 \(\mathrm{d}_1:F_1\to F_0\)。于是 \(\operatorname{im}\mathrm{d}_1=\ker\varepsilon\)。

	重复此过程：已构造 \(\mathrm{d}_i\) 后，选自由模 \(F_{i+1}\) 满射到 \(\ker \mathrm{d}_i\)，便得到 \(\mathrm{d}_{i+1}\)。每一步都按核构造，所以相邻微分复合为零，且所得列正合。

	若 \(A\) Noether 且 \(M\) 有限生成，先取有限秩 \(F_0\)。由\cref{g5:prop:noetherian-finite-modules}，它的子模 \(\ker\varepsilon\) 有限生成，故 \(F_1\) 可取有限秩；逐次重复同一论证即可。这里得到的是每一项秩有限，消解的非零项仍可能有无限多个。
\end{proof}

例如，若 \(f\in A\) 不是零因子，则
\[
	0\longrightarrow A\xrightarrow{\,f\,}A
	\longrightarrow A/(f)\longrightarrow0
\]
就是自由消解。左端的单射正是非零因子条件。若 \(f\) 是零因子，这个序列在左边留下 \((0:_A f)\)，必须继续用自由模覆盖这个核，不能直接把两项呈示当作完整消解。

\ProseParagraphBreak

本节也会使用投射消解。根据\cref{g5:def:projective-module}，投射性就是沿满射提升映射的性质。任意秩的自由模都具有这一性质，因为可以逐个提升基向量的像。进一步，任意投射模 \(P\) 都是某个自由模的直和项：取满射 \(F\to P\)，将 \(\operatorname{id}_P\) 提升得到分裂，便有 \(F\cong P\oplus K\)；反向结论由限制自由模的提升得到。因而投射模平坦：自由模张量对应直和，保留单射，直和项也保留这一性质。此外，\(\operatorname{Hom}_A(P,-)\) 正合，因为其左正合性已由\cref{g5:prop:hom-left-exact}证明，而右端的满性恰是投射性的提升要求。这些说明也适用于无限秩情形。自由与投射消解的构造可对照\cite[附录B，第277--278页]{Matsumura1989}。

\subsection{比较定理和链同伦}
\label{g5:subsec:0801:02}

自由消解依赖生成元的选择。同一个模可以具有大小很不相同的消解，因此用消解计算出的对象，必须先证明不依赖这些选择。所需的比较不要求两份消解逐项同构，而是把模同态提升到各阶，并说明不同提升只相差一种在同调中消失的误差。

\begin{definition}\label{g5:def:chain-homotopy}
	设 \(C_\bullet,D_\bullet\) 为交换含幺环 \(A\) 上的链复形。若同态族 \(u_i:C_i\to D_i\) 满足
	\[
		\mathrm{d}_i^D u_i=u_{i-1}\mathrm{d}_i^C,
	\]
	则称 \(u=(u_i):C_\bullet\to D_\bullet\) 为\term[lianyingshe]{链映射}[chain map]。对两张链映射 \(u,v\)，若同态族 \(h_i:C_i\to D_{i+1}\) 满足
	\[
		u_i-v_i=\mathrm{d}_{i+1}^D h_i+h_{i-1}\mathrm{d}_i^C,
	\]
	则称 \(h=(h_i)\) 为 \(u,v\) 之间的一个\term[liantonglun]{链同伦}[chain homotopy]，记 \(u\simeq v\)。对非负次数复形，约定 \(h_{-1}=0\)。
\end{definition}

链映射将循环送到循环、边界送到边界，因此诱导同调映射。链同伦则使诱导映射相同：若 \(\mathrm{d}_i^Cz=0\)，则
\[
	u_i(z)-v_i(z)=\mathrm{d}_{i+1}^Dh_i(z)
\]
是边界。若两复形之间有链映射 \(u,v\)，使 \(vu\simeq1\)、\(uv\simeq1\)，就称它们\term[liantonglundengjia]{链同伦等价}[chain homotopy equivalence]；此时全部同调模自然同构。下面的\term[xiaojiebijiaodingli]{消解比较定理}[comparison theorem for resolutions]使这种同伦比较适用于任意模同态。

\begin{theorem}[消解比较定理]\label{g5:thm:resolution-comparison}
	设 \(A\) 为交换含幺环，\(P_\bullet\to M\)、\(Q_\bullet\to N\) 为投射消解，\(f:M\to N\) 为同态。则 \(f\) 可以提升为链映射 \(u:P_\bullet\to Q_\bullet\)，使
	\(\varepsilon_Q u_0=f\varepsilon_P\)。任意两张这样的提升链同伦。特别地，同一个模的任意两份投射消解链同伦等价。
\end{theorem}
\begin{proof}
	先构造提升。由于 \(\varepsilon_Q:Q_0\to N\) 满射，\(P_0\) 的投射性给出 \(u_0\)。若已构造低于 \(i\) 次的映射，则 \(u_{i-1}\mathrm{d}_i^P\) 的像落在 \(\ker \mathrm{d}_{i-1}^Q\) 中；在 \(i=1\) 时，这里的核指 \(\ker\varepsilon_Q\)。目标消解正合，故
	\[
		Q_i\twoheadrightarrow\ker \mathrm{d}_{i-1}^Q.
	\]
	用 \(P_i\) 的投射性提升，得到满足 \(\mathrm{d}_i^Qu_i=u_{i-1}\mathrm{d}_i^P\) 的 \(u_i\)。于是逐阶补成交换图
	\[
		\begin{tikzcd}[column sep=large]
			\cdots\arrow[r]&P_1\arrow[r,"\mathrm{d}_1^P"]\arrow[d,"u_1"]
			&P_0\arrow[r,"\varepsilon_P"]\arrow[d,"u_0"]
			&M\arrow[r]\arrow[d,"f"]&0\\
			\cdots\arrow[r]&Q_1\arrow[r,"\mathrm{d}_1^Q"]
			&Q_0\arrow[r,"\varepsilon_Q"]&N\arrow[r]&0.
		\end{tikzcd}
	\]

	再比较两个提升 \(u,v\)。因 \(\varepsilon_Q(u_0-v_0)=0\)，可以沿 \(Q_1\twoheadrightarrow\ker\varepsilon_Q\) 提升 \(u_0-v_0\)，得到 \(h_0\)，使 \(u_0-v_0=\mathrm{d}_1^Qh_0\)。假设已构造 \(h_0,\ldots,h_{i-1}\)，考虑
	\[
		r_i=u_i-v_i-h_{i-1}\mathrm{d}_i^P.
	\]
	链映射等式和上一阶同伦等式给出
	\[
		\mathrm{d}_i^Qr_i
		=(u_{i-1}-v_{i-1}-\mathrm{d}_i^Qh_{i-1})\mathrm{d}_i^P
		=h_{i-2}\mathrm{d}_{i-1}^P\mathrm{d}_i^P=0.
	\]
	故 \(r_i\) 可由投射性提升到 \(Q_{i+1}\)，即 \(r_i=\mathrm{d}_{i+1}^Qh_i\)。这正是所需同伦。

	最后，对同一个模的两份消解，分别提升恒等映射得到 \(u,v\)。复合 \(vu,uv\) 仍提升恒等映射，因而分别与各消解的恒等链映射同伦。
\end{proof}

比较定理中，正合性保证“需要提升的像确实落在某个核中”，投射性保证“能从下一项选出整个线性映射的提升”。两项条件各有明确的作用；该结论及同伦约定见\cite[附录B，第274、278页]{Matsumura1989}。

\ProseParagraphBreak

张量及 \(\operatorname{Hom}\) 都保留上述同伦信息。若 \(u-v=\mathrm{d}h+h\mathrm{d}\)，张量任意 \(A\) 模 \(N\) 后，同一个等式变成
\[
	(u-v)\otimes1=\mathrm{d}(h\otimes1)+(h\otimes1)\mathrm{d}.
\]
对 \(\operatorname{Hom}_A(-,N)\)，取上链微分
\[
	\delta^i(\varphi)=\varphi \mathrm{d}_{i+1}.
\]
将 \(h\) 变成降次数映射 \(s^i(\varphi)=\varphi h_{i-1}\)，则预合成给出的两张上链映射满足 \(u^*-v^*=\delta s+s\delta\)。因此两种运算之后，链同伦等价仍诱导同调或上同调的同构。即使张量破坏了原消解的正合性，这些新出现的同调也与所选消解无关。

\subsection{Ext、Tor 与长正合列}
\label{g5:subsec:0801:03}

要将短正合列中的三个模联系起来，先处理复形层面的基本问题：逐项正合的一列复形，在同调上为什么会多出一张跨次数的连接映射。下面直接构造这张映射，并核对各处的正合性。

\begin{proposition}\label{g5:prop:complex-long-exact}
	设 \(A\) 为交换含幺环，且
	\[
		0\longrightarrow C'_\bullet\xrightarrow{\iota}C_\bullet
		\xrightarrow{p}C''_\bullet\longrightarrow0
	\]
	是逐次短正合的链复形列。则存在自然的 \(A\) 线性连接映射
	\(\partial_i:H_i(C'')\to H_{i-1}(C')\)，使
	\[
		\cdots\longrightarrow H_i(C')\longrightarrow H_i(C)
		\longrightarrow H_i(C'')\xrightarrow{\partial_i}
		H_{i-1}(C')\longrightarrow\cdots
	\]
	正合。对上链复形，相应连接映射为
	\(H^i(C'')\to H^{i+1}(C')\)，并有同样的长正合列。这里自然性指：短正合复形列之间的交换图，诱导与所有连接映射交换的同调图。
\end{proposition}
\begin{proof}
	给定循环 \(z''\in C''_i\)，取提升 \(z\in C_i\)。由于 \(p(dz)=\mathrm{d} z''=0\)，存在唯一的 \(y\in C'_{i-1}\) 使 \(dz=\iota(y)\)。又 \(\iota(dy)=\mathrm{d}^2z=0\)，由单射性得 \(dy=0\)。定义
	\[
		\partial_i[z'']=[y].
	\]
	若提升改成 \(z+\iota(w)\)，则 \(y\) 改成 \(y+dw\)，同调类不变。若代表改成 \(z''+dv''\)，取 \(v''\) 的提升 \(v\)，便可把 \(z\) 改成 \(z+dv\)，其微分不变。因此连接映射良定义；对提升作加法或标量乘法，立即得到其 \(A\) 线性。

	在 \(H_i(C)\) 处，若循环 \(z\) 的像在 \(C''\) 中是边界，写 \(p(z)=dv''\)，提升 \(v''\) 为 \(v\)。于是 \(z-dv=\iota(z')\)，且 \(z'\) 是循环，说明该类来自 \(H_i(C')\)。反向包含由 \(p\iota=0\) 给出。

	在 \(H_i(C'')\) 处，若 \(\partial_i[z'']=0\)，则刚才的 \(y=dw\)。把提升改为 \(z-\iota(w)\)，其微分为零，故 \([z'']\) 来自 \(H_i(C)\)。反过来，若能选取循环提升，连接映射当然为零。

	在 \(H_{i-1}(C')\) 处，若循环 \(y\) 满足 \(\iota(y)=dz\)，则 \(p(z)\) 为循环，而且 \(\partial_i[p(z)]=[y]\)。反向包含来自连接映射构造中的同一等式。这三项证明所有位置的正合性。

	对于短正合复形列之间的交换图，把选定的 \(z'',z,y\) 送到另一列，它们仍满足同样的提升与微分等式，所以连接映射交换。这证明自然性。将上链复形按 \(C_i=C^{-i}\) 改写，便得到升次数的版本。
\end{proof}

这就是\term[changzhenghetongdiaolie]{同调长正合列}[long exact homology sequence]。连接映射的意义很具体：右侧循环的任一提升可能不是循环，其微分落入左侧复形，所产生的同调类正好测量能否把提升修正为循环。构造可对照\cite[附录B，第274--275页]{Matsumura1989}。

\begin{definition}\label{g5:def:ext-tor}
	设 \(A\) 为交换含幺环，\(M,N\) 为 \(A\) 模，取 \(M\) 的自由消解 \(P_\bullet\to M\)。对 \(i\geq0\)，定义
	\[
		\operatorname{Tor}_i^A(M,N)
		=H_i(P_\bullet\otimes_A N),\qquad
		\operatorname{Ext}^i_A(M,N)
		=H^i\bigl(\operatorname{Hom}_A(P_\bullet,N)\bigr).
	\]
	其中张量复形的微分为 \(\mathrm{d}_i\otimes1\)，而
	\(\operatorname{Hom}\) 复形的微分为
	\(\delta^i(\varphi)=\varphi \mathrm{d}_{i+1}\)。它们分别给出\term[Torhanzi]{Tor 函子}[Tor functor]与\term[Exthanzi]{Ext 函子}[Ext functor]。
\end{definition}
\symidx{\(\operatorname{Tor}_i^A(M,N)\)：自由消解张量后的第 \(i\) 个同调}
\symidx{\(\operatorname{Ext}^i_A(M,N)\)：自由消解取 \(\operatorname{Hom}\) 后的第 \(i\) 个上同调}

比较定理及其后的同伦计算证明，这一定义不依赖自由消解，也可以用任意投射消解计算。更准确地说，两份消解之间提升恒等映射，给出与提升选择无关的同构；三份消解之间的这些同构满足复合相容性，因为两种复合都提升恒等映射。

\ProseParagraphBreak

同样，模同态 \(f:M\to M'\) 的链提升，诱导
\[
	\operatorname{Tor}_i^A(M,N)\longrightarrow
	\operatorname{Tor}_i^A(M',N),\qquad
	\operatorname{Ext}^i_A(M',N)\longrightarrow
	\operatorname{Ext}^i_A(M,N).
\]
第二变量的同态直接对张量或 \(\operatorname{Hom}\) 逐项作用，两者在这个变量都协变。不同链提升同伦，所以诱导映射唯一；复合与恒等映射由比较定理相容。同态的加法也相容，因为两个链提升之和提升原同态之和。两个变量中的作用彼此交换，且全部映射 \(A\) 线性，这就是所需的自然性。在零次处，张量的右正合性及 \(\operatorname{Hom}\) 的核描述给出
\begin{equation}\label{g5:eq:ext-tor-degree-zero}
	\operatorname{Tor}_0^A(M,N)\cong M\otimes_A N,\qquad
	\operatorname{Ext}^0_A(M,N)\cong\operatorname{Hom}_A(M,N).
\end{equation}
若 \(M\) 投射，可用仅有零次项 \(M\) 的投射消解，故所有正次数的 \(\operatorname{Ext}^i_A(M,N)\) 与 \(\operatorname{Tor}_i^A(M,N)\) 都为零。定义与这些基本性质见\cite[附录B，第278--280页]{Matsumura1989}。

\begin{proposition}\label{g5:prop:ext-annihilator}
	设 \(A\) 为交换含幺环，\(M,N\) 为 \(A\) 模，\(i\geq0\)。在
	\(\operatorname{Ext}^i_A(M,N)\) 与 \(\operatorname{Tor}_i^A(M,N)\) 上，乘以 \(a\in A\) 的自然标量作用，既等于第一变量乘以 \(a\) 所诱导的映射，也等于第二变量乘以 \(a\) 所诱导的映射。因此
	\[
		\bigl(\operatorname{Ann}_A M+\operatorname{Ann}_A N\bigr)
		\operatorname{Ext}^i_A(M,N)=0,
	\]
	对 \(\operatorname{Tor}_i^A(M,N)\) 也有相同结论。特别地，若 \((A,\mathfrak m,k)\) 为局部环，则全部 \(\operatorname{Ext}^i_A(k,N)\) 与 \(\operatorname{Tor}_i^A(k,N)\) 自然为 \(k\) 向量空间。上述函子与有限直和相容；在第二变量按分量的识别下，同一模的有限直和之间由 \(A\) 系数矩阵给出的映射，诱导的映射仍由同一标量矩阵给出。
\end{proposition}
\begin{proof}
	取 \(M\) 的自由消解 \(P_\bullet\)。各项乘以 \(a\) 构成链映射，并提升 \(M\) 上乘以 \(a\) 的映射。对张量与 \(\operatorname{Hom}\)，分别有
	\[
		(ap)\otimes n=p\otimes(an)=a(p\otimes n),\qquad
		\varphi(ap)=a\varphi(p).
	\]
	所以两个变量诱导的作用在复形上就等于自然标量作用。

	若 \(aM=0\)，各项乘以 \(a\) 的链映射与零链映射都提升 \(M\to M\) 的零映射。由比较定理，它们链同伦，故在所计算的同调与上同调上作用为零。若 \(aN=0\)，第二变量的作用已在复形上为零。于是两个零化理想之和也作用为零。对 \(M=k=A/\mathfrak m\)，取其零化理想 \(\mathfrak m\) 即得向量空间结构。

	最后，张量积及 \(\operatorname{Hom}_A(P_i,-)\) 均与有限直和相容，而直和复形的核与像逐分量计算。因此
	\[
		\operatorname{Ext}^i_A(M,N^r)
		\cong\operatorname{Ext}^i_A(M,N)^r,
	\]
	对 \(\operatorname{Tor}\) 也一样。若 \(N^r\to N^s\) 由矩阵 \((a_{uv})\) 给出，逐项作用就是先乘各 \(a_{uv}\) 再相加；由已经证明的标量作用与加法相容性，诱导映射恰为同一矩阵。第一变量的有限直和也逐项分解，\(\operatorname{Ext}\) 在该变量相应使用转置矩阵。
\end{proof}

这个证明不仅说明某个抽象模被 \(\mathfrak m\) 杀死，还确定了具体标量映射所诱导的作用。后面把一个正则元素的乘法放入短正合列时，正需要这种与原映射相容的零化结论。

\ProseParagraphBreak

为处理第一变量的短正合列，需要把两端消解拼成中间消解。这一构造称为\term[matiyinli]{马蹄引理}[horseshoe lemma]。

\begin{lemma}[马蹄引理]\label{g5:lem:horseshoe}
	设 \(A\) 为交换含幺环，\(0\to M'\to M\to M''\to0\) 为短正合列，且已选 \(M'\)、\(M''\) 的投射消解 \(P'_\bullet,P''_\bullet\)。则存在 \(M\) 的投射消解 \(P_\bullet\) 及增广相容的逐次分裂短正合列
	\[
		0\longrightarrow P'_\bullet\longrightarrow P_\bullet
		\longrightarrow P''_\bullet\longrightarrow0,
		\qquad P_i=P'_i\oplus P''_i.
	\]
	若原来两份消解自由，则中间消解也自由。
\end{lemma}
\begin{proof}
	记原列的映射为 \(\iota:M'\to M\)、\(p:M\to M''\)。由 \(P''_0\) 的投射性，将其增广 \(\varepsilon''\) 提升为 \(s_0:P''_0\to M\)。定义
	\[
		\varepsilon:P'_0\oplus P''_0\longrightarrow M,\qquad
		(u,v)\longmapsto\iota\varepsilon'(u)+s_0(v).
	\]
	此映射满射：先提升一个元素在 \(M''\) 中的像，剩下的差在 \(\iota(M')\) 内，再由 \(P'_0\) 提升。令三个增广的核为 \(K'_1,K_1,K''_1\)。投影到第二分量给出
	\[
		0\longrightarrow K'_1\longrightarrow K_1
		\longrightarrow K''_1\longrightarrow0.
	\]
	核的描述直接来自 \((u,0)\in K_1\iff u\in K'_1\)。为证满射，给定 \(v\in K''_1\)，有 \(s_0(v)\in\iota(M')\)，故可选 \(u\) 使 \(\iota\varepsilon'(u)=-s_0(v)\)，于是 \((u,v)\in K_1\)。

	现在以这条核的短正合列代替原列，并用既有满射 \(P'_1\to K'_1\)、\(P''_1\to K''_1\) 重复构造，得到 \(P_1=P'_1\oplus P''_1\to K_1\)。与 \(K_1\hookrightarrow P_0\) 复合，即得中间微分。继续处理核，就得到全部 \(P_i\) 及微分。每一步的投影、包含都与构造相容，故所得三列构成短正合复形列；各次按直和分裂，而中间列按核的构造正合。
\end{proof}

其中的直和分解只要求在每个次数成立，中间微分可能含有从 \(P''_i\) 到 \(P'_{i-1}\) 的分量，故中间复形一般不是两侧复形的直和。正是这些分量记录原短正合列的相互联系。

\begin{proposition}\label{g5:prop:ext-tor-long-exact}
	设 \(A\) 为交换含幺环。对任意 \(A\) 模，\(\operatorname{Ext}\) 与 \(\operatorname{Tor}\) 有下列自然的 \(A\) 线性长正合列。
	\begin{enumerate}
		\item 对 \(0\to N'\to N\to N''\to0\)，固定第一变量 \(M\)，有
		\[
			\begin{aligned}
				0&\longrightarrow\operatorname{Hom}_A(M,N')
				\longrightarrow\operatorname{Hom}_A(M,N)
				\longrightarrow\operatorname{Hom}_A(M,N'')\\
				&\longrightarrow\operatorname{Ext}^1_A(M,N')
				\longrightarrow\operatorname{Ext}^1_A(M,N)
				\longrightarrow\operatorname{Ext}^1_A(M,N'')
				\\
				&\longrightarrow\operatorname{Ext}^2_A(M,N')
				\longrightarrow\cdots.
			\end{aligned}
		\]
		\item 对 \(0\to M'\to M\to M''\to0\)，固定第二变量 \(N\)，有
		\[
			\begin{aligned}
				0&\longrightarrow\operatorname{Hom}_A(M'',N)
				\longrightarrow\operatorname{Hom}_A(M,N)
				\longrightarrow\operatorname{Hom}_A(M',N)\\
				&\longrightarrow\operatorname{Ext}^1_A(M'',N)
				\longrightarrow\operatorname{Ext}^1_A(M,N)
				\longrightarrow\operatorname{Ext}^1_A(M',N)
				\\
				&\longrightarrow\operatorname{Ext}^2_A(M'',N)
				\longrightarrow\cdots.
			\end{aligned}
		\]
		\item 对第一变量的短正合列，对各 \(i\geq1\) 有
		\[
			\begin{aligned}
				\cdots&\longrightarrow\operatorname{Tor}_i^A(M',N)
				\longrightarrow\operatorname{Tor}_i^A(M,N)
				\longrightarrow\operatorname{Tor}_i^A(M'',N)\\
				&\longrightarrow\operatorname{Tor}_{i-1}^A(M',N)
				\longrightarrow\cdots,
			\end{aligned}
		\]
		末端为
		\[
			\operatorname{Tor}_1^A(M'',N)\longrightarrow M'\otimes_A N
			\longrightarrow M\otimes_A N
			\longrightarrow M''\otimes_A N\longrightarrow0.
		\]
		第二变量也有同样形式的长正合列，其连接映射为
		\(\operatorname{Tor}_i^A(M,N'')\to\operatorname{Tor}_{i-1}^A(M,N')\)。
	\end{enumerate}
\end{proposition}
\begin{proof}
	固定 \(M\) 的投射消解 \(P_\bullet\)。每个 \(P_i\) 投射，故
	\[
		0\to\operatorname{Hom}_A(P_\bullet,N')
		\to\operatorname{Hom}_A(P_\bullet,N)
		\to\operatorname{Hom}_A(P_\bullet,N'')\to0
	\]
	逐次正合。应用\cref{g5:prop:complex-long-exact}的上链版本及零次识别，得到第一列。又每个 \(P_i\) 平坦，所以张量 \(P_\bullet\) 后也有逐次短正合的链复形列，得到 \(\operatorname{Tor}\) 对第二变量的长正合列。

	对第一变量，取马蹄引理中的短正合消解列。它逐次分裂，故分别取 \(\operatorname{Hom}_A(-,N)\) 或张量 \(N\) 后仍逐次正合；前者反转三项次序。再用复形长正合列，就得到第二列及 \(\operatorname{Tor}\) 对第一变量的版本。

	还需核对后一构造的自然性与选择无关性。给定两条模短正合列之间的交换图，可将它提升为马蹄消解列之间的交换图。具体地，先提升两端映射；零次中间映射在两端上的分量已经确定，增广误差落在目标的左端模中，用 \(P''_0\) 的投射性沿左端增广提升这个误差，便补足一个从 \(P''_0\) 到目标左端零次项的分量。在下一次，对已经得到的核的短正合列重复同一过程。这样逐阶得到与包含及投影交换的链映射。复形连接映射的自然性使所得长正合图交换，而比较定理保证每项同调上的映射正是已定义的模同态所诱导的映射。特别地，对原模短正合列的恒等图作此构造，说明不同马蹄选择给出的连接映射在消解的规范识别下相同。
\end{proof}

长正合列把第二章的左正合、右正合性质接续为一条完整的列。例如 \(\operatorname{Hom}_A(M,N)\to\operatorname{Hom}_A(M,N'')\) 中一个映射不能提升，其障碍落在 \(\operatorname{Ext}^1_A(M,N')\)；只有该障碍为零时，正合性才保证提升存在。

\ProseParagraphBreak

定义 \(\operatorname{Tor}\) 时只消解了第一变量，实际计算却常需要选择较容易消解的那一项。下面证明两种计算一致。证明只处理第一象限，因此每个总次数只有有限多个分量。

\begin{lemma}\label{g5:lem:double-complex-collapse}
	设 \(A\) 为交换含幺环，\(D_{p,q}\) 为 \(p,q\geq0\) 的 \(A\) 模，配有两个微分
	\[
		h:D_{p,q}\to D_{p-1,q},\qquad
		v:D_{p,q}\to D_{p,q-1},
		\qquad h^2=v^2=0,\quad hv=vh.
	\]
	负指标处取零。令
	\[
		T_n=\bigoplus_{p+q=n}D_{p,q},\qquad
		\mathrm{d}|_{D_{p,q}}=h+(-1)^p v.
	\]
	若每一列关于 \(v\) 的正次数同调为零，令
	\(X_p=D_{p,0}/vD_{p,1}\)，以 \(h\) 诱导微分，则投影 \(T_\bullet\to X_\bullet\) 诱导
	\[
		H_n(T_\bullet)\cong H_n(X_\bullet)
	\]
	对全部 \(n\geq0\) 成立，且同构对保持两个微分的映射自然。
\end{lemma}
\begin{proof}
	在 \(D_{p,q}\) 上展开 \(\mathrm{d}^2\)，两个交叉项的符号分别为 \((-1)^{p-1}\) 与 \((-1)^p\)，由 \(hv=vh\) 抵消，其余两项为零。因此 \(T_\bullet\) 是链复形。投影只取总次数 \(n\) 中的 \((n,0)\) 分量再取商；落到零行的竖直微分在商中为零，所以这是逐次满射的链映射。

	设其核复形为 \(L_\bullet\)。我们逐个消去分量，证明每个 \(L_n\) 中的循环 \(z\) 都是边界。首先，\(z\) 的 \((n,0)\) 分量属于 \(vD_{n,1}\)。选 \(w_{n,1}\)，使 \((-1)^n v w_{n,1}\) 等于这个分量，则 \(z-\mathrm{d}w_{n,1}\) 的 \((n,0)\) 分量为零。这个修正仍属于 \(L_\bullet\)，因为 \(w_{n,1}\) 的投影为零。

	现在从 \(p=n-1\) 逐次下降到 \(0\)。设 \(p\) 右边的所有列已消去，当前可能非零的最右分量为 \(z_{p,n-p}\)。总微分为零，且右边的列已为零，故
	\[
		vz_{p,n-p}=0.
	\]
	这里 \(n-p>0\)，由该列的正次数同调消失，可以选
	\(w_{p,n-p+1}\)，使 \((-1)^p v w_{p,n-p+1}=z_{p,n-p}\)。减去其总边界，便消去这一分量，只可能在左一列产生新分量。有限次之后全部分量消去，所以 \(z\) 是这些 \(w\) 之和的边界。\(n=0\) 时仅需第一步。

	于是 \(H_n(L)=0\) 对所有 \(n\) 成立。对短正合复形列 \(0\to L\to T\to X\to0\) 使用\cref{g5:prop:complex-long-exact}，即得所需同构。投影本身与给定映射交换，故同构自然。
\end{proof}

这种带有两个方向微分的对象称为\term[shuangfuxing]{双复形}[double complex]，所得 \(T_\bullet\) 称为其\term[zongfuxing]{总复形}[total complex]。引理的关键是固定总次数中的列数有限，故逐列消去确实终止；这一特例见\cite[附录B，定理B1，第275--276页]{Matsumura1989}。

\begin{proposition}\label{g5:prop:tor-balanced}
	设 \(A\) 为交换含幺环，\(M,N\) 为 \(A\) 模，\(P_\bullet\to M\)、\(Q_\bullet\to N\) 为投射消解。则有自然同构
	\[
		\operatorname{Tor}_i^A(M,N)
		\cong H_i(M\otimes_A Q_\bullet)
		\cong\operatorname{Tor}_i^A(N,M)
		\qquad(i\geq0).
	\]
	因此可以消解任意一个变量来计算 \(\operatorname{Tor}\)。
\end{proposition}
\begin{proof}
	取 \(D_{p,q}=P_p\otimes_A Q_q\)，两个微分分别为 \(\mathrm{d}^P\otimes1\) 与 \(1\otimes \mathrm{d}^Q\)，它们交换。因为 \(P_p\) 平坦，第 \(p\) 列的正次数同调为零，零次同调为 \(P_p\otimes_A N\)。由上一引理，其总复形 \(T\) 满足
	\[
		H_i(T)\cong H_i(P_\bullet\otimes_A N).
	\]
	同样，每个 \(Q_q\) 平坦，所以各行的正次数同调为零，零次同调为 \(M\otimes_A Q_q\)。也可以沿行使用引理：交换两个指标时，对 \((p,q)\) 分量乘以 \((-1)^{pq}\)，便把两种总微分相互对应。这给出
	\[
		H_i(T)\cong H_i(M\otimes_A Q_\bullet).
	\]
	两者合成即为第一同构。交换纯张量的两个因子，再按定义计算，得到第二同构。双复形的投影及交换映射对链映射自然；比较定理保证这些识别也与消解选择相容。
\end{proof}

\subsection{维数平移和合冲模}
\label{g5:subsec:0801:04}

一次自由呈示的核是原模的关系模；继续取核，就把高阶关系转成另一个模的低阶关系。长正合列使这种转移成为精确的同构，而不仅是对消解图形的描述。

\begin{definition}\label{g5:def:syzygy-module}
	设 \(A\) 为交换含幺环，\(M\) 为 \(A\) 模，选定它的投射消解 \(P_\bullet\to M\)。令
	\[
		\Omega^0M=M,\qquad
		\Omega^1M=\ker(P_0\to M),\qquad
		\Omega^rM=\ker(P_{r-1}\to P_{r-2})\quad(r\geq2).
	\]
	称 \(\Omega^rM\) 为该消解的第 \(r\) 个\term[hechongmo]{合冲模}[syzygy module]。消解正合性给出短正合列
	\[
		0\longrightarrow\Omega^{r+1}M\longrightarrow P_r
		\longrightarrow\Omega^rM\longrightarrow0
		\qquad(r\geq0).
	\]
\end{definition}
\symidx{\(\Omega^rM\)：选定消解的第 \(r\) 个合冲模}

合冲模依赖消解的选择，记号中省略的是这份已经选定的消解。它们不能直接作为“唯一确定的核”使用。例如给零次自由模增添一个多余生成元时，第一次合冲模通常也会增大。下面的同构说明，正次数的同调障碍仍然受原模控制。

\begin{proposition}\label{g5:prop:dimension-shifting}
	设 \(A\) 为交换含幺环，\(M,N\) 为 \(A\) 模，并选定 \(M\) 的投射消解及其合冲模。则对 \(r\geq0\)、\(i\geq1\)，有自然同构
	\[
		\operatorname{Ext}^i_A(\Omega^rM,N)
		\cong\operatorname{Ext}^{i+r}_A(M,N),\qquad
		\operatorname{Tor}_i^A(\Omega^rM,N)
		\cong\operatorname{Tor}_{i+r}^A(M,N).
	\]
	此外，记 \(K=\Omega^1M\)，有正合列
	\[
		0\longrightarrow\operatorname{Tor}_1^A(M,N)
		\longrightarrow K\otimes_A N
		\longrightarrow P_0\otimes_A N
		\longrightarrow M\otimes_A N\longrightarrow0
	\]
	以及
	\[
		\begin{aligned}
			0&\longrightarrow\operatorname{Hom}_A(M,N)
			\longrightarrow\operatorname{Hom}_A(P_0,N)
			\longrightarrow\operatorname{Hom}_A(K,N)\\
			&\longrightarrow\operatorname{Ext}^1_A(M,N)
			\longrightarrow0.
		\end{aligned}
	\]
\end{proposition}
\begin{proof}
	将\cref{g5:prop:ext-tor-long-exact}用于
	\(0\to\Omega^{r+1}M\to P_r\to\Omega^rM\to0\)。因为 \(P_r\) 投射，包含它的正次数 \(\operatorname{Ext}\) 与 \(\operatorname{Tor}\) 都消失，所以长正合列中留下同构
	\[
		\operatorname{Ext}^i_A(\Omega^{r+1}M,N)
		\cong\operatorname{Ext}^{i+1}_A(\Omega^rM,N),
	\]
	\[
		\operatorname{Tor}_{i+1}^A(\Omega^rM,N)
		\cong\operatorname{Tor}_i^A(\Omega^{r+1}M,N)
		\qquad(i\geq1).
	\]
	逐次复合便得一般结论。取 \(r=0\) 并保留零次项，得到所述两条有限正合列。自然性继承自长正合列，既适用于第二变量，也适用于选定消解之间诱导的合冲模映射。
\end{proof}

这些同构称为\term[weishupingyi]{维数平移}[dimension shifting]；这里移动的是同调次数，而不是环的 Krull 维数。尤其
\[
	\operatorname{Ext}^1_A(M,N)
	\cong
	\operatorname{Hom}_A(K,N)/
	\operatorname{im}\bigl(\operatorname{Hom}_A(P_0,N)\bigr)
\]
说明，关系模上的同态能否延拓到零次自由模，其障碍恰由一次 \(\operatorname{Ext}\) 记录。相应的合冲模与次数平移见\cite[附录B，第280页]{Matsumura1989}。

\begin{proposition}\label{g5:prop:syzygy-stable}
	设 \(A\) 为交换含幺环，且
	\[
		0\to K\to P\xrightarrow{\pi}M\to0,\qquad
		0\to K'\to P'\xrightarrow{\pi'}M\to0
	\]
	正合，其中 \(P,P'\) 投射。则
	\[
		K\oplus P'\cong K'\oplus P.
	\]
\end{proposition}
\begin{proof}
	在 \(P\oplus P'\) 中取子模
	\[
		T=\{(p,p'):\pi(p)=\pi'(p')\}.
	\]
	投影到 \(P'\) 满射，因为 \(\pi\) 满射；其核为 \(K\)。故
	\(0\to K\to T\to P'\to0\) 由 \(P'\) 的投射性分裂，得到 \(T\cong K\oplus P'\)。交换两份呈示，同理得 \(T\cong K'\oplus P\)。两者合成即得所述同构。
\end{proof}

所以不同呈示的关系模虽未必同构，补上相应的投射直和项后却一致。另一方面，是否需要继续增加关系，还有一个完全由 \(\operatorname{Ext}\) 表述的判据。

\begin{proposition}\label{g5:prop:projectivity-ext}
	设 \(A\) 为交换含幺环，\(M\) 为 \(A\) 模。则 \(M\) 投射，当且仅当对每个 \(A\) 模 \(N\)，都有
	\[
		\operatorname{Ext}^1_A(M,N)=0.
	\]
	若 \(M\) 有限呈示，这还等价于：对每个素理想 \(\mathfrak p\)，局部化 \(M_{\mathfrak p}\) 是有限秩自由 \(A_{\mathfrak p}\) 模。
\end{proposition}
\begin{proof}
	投射模的正次数 \(\operatorname{Ext}\) 已知为零。反过来，取自由满射 \(F\to M\)，核为 \(K\)。对短正合列 \(0\to K\to F\to M\to0\) 在第二变量取 \(\operatorname{Ext}_A(M,-)\)，得到
	\[
		\operatorname{Hom}_A(M,F)\longrightarrow
		\operatorname{Hom}_A(M,M)\longrightarrow
		\operatorname{Ext}^1_A(M,K)=0.
	\]
	因此 \(\operatorname{id}_M\) 提升为 \(M\to F\)，使自由满射分裂。\(M\) 是自由模直和项，故投射。有限呈示情形的最后一项正是\cref{g5:thm:finite-projective-local}：有限呈示使局部的自由基和分裂能够在有限主开覆盖上拼合。
\end{proof}

量词“对每个 \(N\)”不能由只检验一个任意指定的模替代。这个判据在合冲模上应用一次，就给出有限长度消解的检测。

\begin{proposition}\label{g5:prop:resolution-length-ext}
	设 \(A\) 为交换含幺环，\(M\) 为 \(A\) 模，\(r\geq0\) 为整数。下列条件等价：
	\begin{equivalences}
		\item 存在正合列 \(0\to P_r\to\cdots\to P_0\to M\to0\)，其中每个 \(P_j\) 投射；
		\item 对每个 \(A\) 模 \(N\)，有 \(\operatorname{Ext}^{r+1}_A(M,N)=0\)；
		\item \(M\) 的任意投射消解中的 \(\Omega^rM\) 投射。
	\end{equivalences}
	若这些条件成立，则对每个 \(N\) 及每个 \(i>r\)，均有
	\(\operatorname{Ext}^i_A(M,N)=0\) 和 \(\operatorname{Tor}_i^A(M,N)=0\)。
\end{proposition}
\begin{proof}
	第一条件下，用这份有限投射消解直接计算，超过 \(r\) 次的复形项为零，所以相应同调与上同调均为零，特别得到第二条件。

	由维数平移，对任意选定的消解有
	\[
		\operatorname{Ext}^1_A(\Omega^rM,N)
		\cong\operatorname{Ext}^{r+1}_A(M,N).
	\]
	若第二条件成立，\cref{g5:prop:projectivity-ext}使 \(\Omega^rM\) 投射，得到第三条件。最后，若 \(r\geq1\)，截断消解为
	\[
		0\longrightarrow\Omega^rM\longrightarrow P_{r-1}
		\longrightarrow\cdots\longrightarrow P_0
		\longrightarrow M\longrightarrow0,
	\]
	就是第一条件所需的投射消解。\(r=0\) 时第三条件表示 \(M\) 自身投射，取零次消解即可。
\end{proof}

\begin{proposition}\label{g5:prop:flatness-tor}
	设 \(A\) 为交换含幺环，\(M\) 为 \(A\) 模。下列条件等价：
	\begin{equivalences}
		\item \(M\) 平坦；
		\item 对每个 \(A\) 模 \(N\) 及每个 \(i\geq1\)，有 \(\operatorname{Tor}_i^A(M,N)=0\)；
		\item 对每个 \(A\) 模 \(N\)，有 \(\operatorname{Tor}_1^A(M,N)=0\)；
		\item 对每个有限生成理想 \(I\subseteq A\)，有 \(\operatorname{Tor}_1^A(A/I,M)=0\)。
	\end{equivalences}
\end{proposition}
\begin{proof}
	若 \(M\) 平坦，取 \(N\) 的自由消解 \(Q_\bullet\)。其各段短正合列张量 \(M\) 后仍正合，因此 \(M\otimes_A Q_\bullet\) 的正次数同调都为零。\cref{g5:prop:tor-balanced}给出第二条件，它当然蕴含第三条件。

	若第三条件成立，对任意单射 \(L\hookrightarrow E\) 取商 \(N=E/L\)。第二变量的 \(\operatorname{Tor}\) 长正合列包含
	\[
		0=\operatorname{Tor}_1^A(M,N)\longrightarrow M\otimes_A L
		\longrightarrow M\otimes_A E.
	\]
	故张量 \(M\) 保留任意单射，即 \(M\) 平坦。

	最后，对 \(0\to I\to A\to A/I\to0\) 应用第一变量的长正合列。由于 \(A\) 投射，得到
	\[
		0\longrightarrow\operatorname{Tor}_1^A(A/I,M)
		\longrightarrow I\otimes_A M\longrightarrow M.
	\]
	所以第四条件恰好表示每个有限生成理想的乘法映射 \(I\otimes_A M\to M\) 单射；由\cref{g5:thm:flat-ideal-criterion}，这等价于平坦性。
\end{proof}

这个判据把第六章的理想张量判据改写成同调语言。一旦找到一个非零的一次 \(\operatorname{Tor}\)，就找到了张量失去单射性的实际核；而投射性的对应障碍，则是 \(\operatorname{Ext}^1\) 中不能提升或延拓的映射。

\subsection{非平坦张量的显式 Tor}
\label{g5:subsec:0801:05}

计算时先选择最短或最容易写出的消解，再写出张量后的微分。不能只看消解中自由模的秩：原来单射的乘法映射，在新的模上可能成为零映射，也可能仍是单射。

\begin{proposition}\label{g5:prop:principal-quotient-ext-tor}
	设 \(A\) 为交换含幺环，\(f\in A\) 不是零因子，\(N\) 为 \(A\) 模。则
	\[
		\begin{aligned}
			\operatorname{Tor}_1^A(A/(f),N)&\cong(0:_N f),&
			\operatorname{Tor}_0^A(A/(f),N)&\cong N/fN,\\
			\operatorname{Ext}^0_A(A/(f),N)&\cong(0:_N f),&
			\operatorname{Ext}^1_A(A/(f),N)&\cong N/fN.
		\end{aligned}
	\]
	当 \(i\geq2\) 时，\(\operatorname{Tor}_i^A(A/(f),N)\) 与
	\(\operatorname{Ext}^i_A(A/(f),N)\) 都为零。
\end{proposition}
\begin{proof}
	使用自由消解 \(0\to A\xrightarrow{f}A\to A/(f)\to0\)。张量 \(N\) 得到链复形
	\[
		0\longrightarrow N\xrightarrow{\,f\,}N\longrightarrow0,
	\]
	其中左边位于次数一，故一次同调为乘以 \(f\) 的核，零次同调为其余核。取 \(\operatorname{Hom}_A(-,N)\) 得到同样的乘法映射，但左边位于上链次数零、右边位于次数一。因此零次上同调为核、一次上同调为余核。其他次数没有非零项，结论随之成立。
\end{proof}

例如取 \(A=\mathbb Z\)、\(m,n\geq1\)。令 \(g=\gcd(m,n)\)，则
\[
	\operatorname{Tor}_1^{\mathbb Z}(\mathbb Z/m,\mathbb Z/n)
	\cong\mathbb Z/g,\qquad
	\operatorname{Ext}^1_{\mathbb Z}(\mathbb Z/m,\mathbb Z/n)
	\cong\mathbb Z/g.
\]
第一个同构来自 \(\mathbb Z/n\) 中被 \(m\) 杀死的元素：它们由 \(n/g\) 的类生成，且这个生成元的阶为 \(g\)。第二个同构来自商
\((\mathbb Z/n)/m(\mathbb Z/n)\)。当 \(m=n>1\) 时，乘以 \(m\) 的单射 \(\mathbb Z\to\mathbb Z\) 张量 \(\mathbb Z/m\) 后成为零映射，而相应的一次 \(\operatorname{Tor}\) 正好是留下的整个核。

\begin{proposition}\label{g5:prop:tor-ideal-intersection}
	设 \(A\) 为交换含幺环，\(I,J\subseteq A\) 为任意理想。则有自然同构
	\[
		\operatorname{Tor}_1^A(A/I,A/J)
		\cong (I\cap J)/IJ.
	\]
\end{proposition}
\begin{proof}
	由 \(0\to I\to A\to A/I\to0\) 的长正合列及 \(\operatorname{Tor}_1^A(A,A/J)=0\)，有
	\[
		0\longrightarrow\operatorname{Tor}_1^A(A/I,A/J)
		\longrightarrow I\otimes_A A/J
		\longrightarrow A\otimes_A A/J.
	\]
	将右边两项识别为 \(I/IJ\) 与 \(A/J\)，中间映射就是
	\(a+IJ\mapsto a+J\)。它的核由 \(a\in I\cap J\) 给出，故为 \((I\cap J)/IJ\)。
\end{proof}

因此一次 \(\operatorname{Tor}\) 比较的是理想的交与乘积；即使两个闭子集的集合交相同，相应商模的关系仍可能不同。这个公式没有对理想添加有限生成条件。后面的几何模型将把它与具体交的方程结合起来。

\begin{example}\label{g5:exm:dual-numbers-infinite-tor}
	设 \(k\) 为任意域，\(A=k[\varepsilon]/(\varepsilon^2)\)，并把
	\(k=A/(\varepsilon)\) 看作 \(A\) 模。则
	\[
		\cdots\xrightarrow{\varepsilon}A
		\xrightarrow{\varepsilon}A
		\xrightarrow{\varepsilon}A\longrightarrow k\longrightarrow0
	\]
	是自由消解。确实，\(A\) 中每个元素唯一写为 \(a+b\varepsilon\)，乘以 \(\varepsilon\) 得 \(a\varepsilon\)，所以每个中间位置的核与像都为 \((\varepsilon)\)，增广的核也为该理想。

	张量 \(k\) 后，全部微分为零。因此
	\[
		\operatorname{Tor}_i^A(k,k)\cong k\qquad(i\geq0).
	\]
	同样，取 \(\operatorname{Hom}_A(-,k)\) 后全部上链微分为零，得到
	\[
		\operatorname{Ext}^i_A(k,k)\cong k\qquad(i\geq0).
	\]
	这些结论对任意特征成立。它们证明 \(k\) 不平坦，而且不存在有限长度的投射消解：若有长度不超过 \(r\) 的消解，超过 \(r\) 次的同调就应全部消失。

	一次障碍也能从一个短正合列直接看到。因为 \((\varepsilon)\cong k\) 作为 \(A\) 模，包含 \((\varepsilon)\hookrightarrow A\) 张量 \(k\) 后变成
	\[
		(\varepsilon)\otimes_A k\cong k
		\longrightarrow A\otimes_A k\cong k,
	\]
	这张映射把 \(\varepsilon\otimes1\) 送到零。长正合列将其非零核识别为
	\(\operatorname{Tor}_1^A(k,k)\)。无限消解进一步表明，同样的关系在每个更高次数继续出现。
\end{example}

自由消解把“生成元、关系及高阶关系”组织成可以比较的复形；链同伦保证其同调不随表示改变。长正合列与维数平移随后把障碍从一项传到另一项。下一节的 Koszul 复形将给出一类由有限个元素直接写出的自由复形，并用正则性判断它何时就是所需消解。
% 纲要条目排为编号小节标题，并进入目录与书签；说明段仍保留于纲要注释中。

% !TeX root = ../../Book5.tex
% 纲要标题：### 8.2 正则序列和 Koszul
% 纲要位置：几何学计划纲要.md，第 5311 行。
% 正文状态：初稿；尚未正式审核。

% 纲要内容开始
% 
% * 逐次非零因子及末端非零条件
% * 外代数上的微分
% * 正则序列给出消解
% * 局部有限模上的反向正合判据
% * 交换次序和基变换的限制
% 
% 纲要内容结束

\section{正则序列和 Koszul}
\label{g5:sec:0802}

一个方程 \(x\) 在模 \(M\) 上没有零因子，意味着乘法映射 \(M\xrightarrow{x}M\) 单射。增加第二个方程时，应该在除去第一个方程后的商模上检查同一条件。Koszul 复形把这串逐次检查合成一个有限复形：一次项记录方程，二次项记录方程间的交换关系，更高项继续记录这些关系之间的相容性。正则序列使正次数上的同调全部消失；在局部 Noether 环的有限模上，反过来只检查第一同调就足够。

\subsection{逐次非零因子及末端非零条件}
\label{g5:subsec:0802:01}

第三章已在环上定义正则序列，见\subsecref{g5:subsec:0304:03}。现在将逐次乘法的条件用于一个指定的模；同一元素可能在这个模上没有零因子，却在底环上有零因子。

\begin{definition}\label{g5:def:module-regular-sequence}
	设 \(A\) 为交换含幺环，\(M\) 为 \(A\) 模。元素 \(x\in A\) 称为 \(M\) 上的非零因子，若乘法映射 \(M\xrightarrow{x}M\) 单射。有限序列 \(x_1,\ldots,x_r\in A\) 称为 \(M\) 正则序列，若：
	\begin{enumerate}
		\item 对每个 \(1\leq j\leq r\)，乘以 \(x_j\) 在
		\[
			M_{j-1}=M/(x_1,\ldots,x_{j-1})M
		\]
		上单射；
		\item 末端商模 \(M_r=M/(x_1,\ldots,x_r)M\) 非零。
	\end{enumerate}
	\(M=A\) 时称为 \(A\) 正则序列，或环上的正则序列。空序列在 \(M\ne0\) 时也按此定义为正则序列。
\end{definition}

末端条件保证每一步仍有非零对象可供下一步研究。它不能只从逐次单射推出。例如在 \(\mathbb Z\) 上，乘以 \(2\) 单射，乘以 \(3\) 在 \(\mathbb Z/2\mathbb Z\) 上也单射，但 \(\mathbb Z/(2,3)=0\)，所以 \(2,3\) 不构成这里的正则序列。定义及末端条件可对照\cite[第16节开头，第123页]{Matsumura1989}。

\begin{proposition}\label{g5:prop:regular-sequence-terminal-quotient}
	设 \((A,\mathfrak m)\) 为交换含幺局部环，\(M\ne0\) 为有限生成 \(A\) 模，且 \(x_1,\ldots,x_r\in\mathfrak m\)。则
	\[
		M/(x_1,\ldots,x_r)M\ne0.
	\]
	因此在这些条件下，正则性只需逐次检查非零因子条件。
\end{proposition}
\begin{proof}
	若末端商为零，则
	\(M=(x_1,\ldots,x_r)M\subseteq\mathfrak mM\subseteq M\)，所以 \(\mathfrak mM=M\)。由\cref{g5:thm:nakayama-local}得到 \(M=0\)，与假设矛盾。
\end{proof}

每个正则序列的前段仍正则，删去前段后留下的序列则在相应商模上正则。这由逐次条件直接得到；各中间商都满射到非零的末端商，故不会提前变为零。在多项式环 \(k[X_1,\ldots,X_n]\) 中，\(X_1,\ldots,X_r\)、\(r\leq n\)，便是基本例子：每次取商后仍是剩余变量的多项式环，下一变量的乘法单射，末端商也非零。

\subsection{外代数上的微分}
\label{g5:subsec:0802:02}

设 \(E=A^r\)，基为 \(e_1,\ldots,e_r\)。其\term[waidaishu]{外代数}[exterior algebra]记为 \(\bigwedge_A E\)。这里可以直接使用自由基描述：\(\bigwedge_A^iE\) 是以
\[
	e_{j_1}\wedge\cdots\wedge e_{j_i},
	\qquad 1\leq j_1<\cdots<j_i\leq r,
\]
为基的自由模，秩为 \(\binom ri\)。规定 \(\bigwedge_A^0E=A\)，并在 \(i<0\) 或 \(i>r\) 时取零。乘法由并列楔积、交换相邻不同基向量改变符号、出现重复基向量时取零来定义，再作多重线性延拓。这些规则给出结合乘法，并满足
\[
	e_i\wedge e_j=-e_j\wedge e_i,
	\qquad e_i\wedge e_i=0.
\]
在特征二时也必须保留第二条规则；仅有第一条不能推出重复项为零。

\begin{definition}\label{g5:def:koszul-complex}
	设 \(A\) 为交换含幺环，\(\mathbf x=(x_1,\ldots,x_r)\in A^r\)，\(M\) 为 \(A\) 模。取 \(E=A^r\)，定义
	\[
		K_i(\mathbf x;M)=\bigwedge_A^iE\otimes_A M.
	\]
	在递增下标的基楔积上，令
	\begin{equation}\label{g5:eq:koszul-differential}
		\begin{aligned}
		&\mathrm d_i(e_{j_1}\wedge\cdots\wedge e_{j_i}\otimes m)\\
		&\qquad=\sum_{s=1}^i(-1)^{s-1}
		 e_{j_1}\wedge\cdots\wedge\widehat{e_{j_s}}
		 \wedge\cdots\wedge e_{j_i}\otimes x_{j_s}m.
		\end{aligned}
	\end{equation}
	其中帽号表示删去该因子，\(\mathrm d_0=0\)。由这些微分构成的链复形称为\term[Koszulfuxing]{Koszul 复形}[Koszul complex]，记为 \(K_\bullet(\mathbf x;M)\)。特别地，\(K_\bullet(\mathbf x;A)\) 的各项均为有限自由 \(A\) 模。
\end{definition}

这个复形以 Jean-Louis Koszul 命名。\footnote{Jean-Louis Koszul（1921--2018），法国数学家，研究 Lie 代数的同调与上同调，并对同调代数的发展作出重要贡献。}微分把一个记录方程的基向量换成相应系数 \(x_j\)，符号则使两种删除顺序相互抵消。

\begin{proposition}\label{g5:prop:koszul-differential-square}
	设 \(A\) 为交换含幺环，\(\mathbf x\in A^r\)，\(M\) 为 \(A\) 模。上面的微分满足 \(\mathrm d_{i-1}\mathrm d_i=0\)。若记
	\[
		H_i(\mathbf x;M)=H_i\bigl(K_\bullet(\mathbf x;M)\bigr),
	\]
	则 \(H_0(\mathbf x;M)=M/(x_1,\ldots,x_r)M\)。
\end{proposition}
\begin{proof}
	在一个基楔积上连续使用两次微分。固定原楔积中第 \(p\)、\(q\) 个因子，\(p<q\)。先删第 \(q\) 个、再删第 \(p\) 个，符号为 \((-1)^{p+q-2}\)；先删第 \(p\) 个以后，原第 \(q\) 个因子的位置变为 \(q-1\)，所以另一顺序的符号为 \((-1)^{p+q-3}\)。两项留下相同的楔积，系数分别为 \(x_{j_p}x_{j_q}\) 与 \(x_{j_q}x_{j_p}\)，由交换性相等，而符号相反。因此每一对删除项都抵消，得到 \(\mathrm d^2=0\)。

	又因 \(K_0=M\)，且 \(\mathrm d_1(e_j\otimes m)=x_jm\)，\(\mathrm d_1\) 的像就是 \(\sum_jx_jM\)，所以零次同调为所述商模。
\end{proof}

例如，单个元素的复形在次数 \(1,0\) 上为
\[
	0\longrightarrow M\xrightarrow{x_1}M\longrightarrow0,
\]
所以 \(H_1(x_1;M)=\{m\in M:x_1m=0\}\)。两个元素时，依次以 \(e_1\wedge e_2\)、\(e_1,e_2\)、\(1\) 作为外代数基，复形为
\[
	0\longrightarrow M
	\xrightarrow{\ m\mapsto(-x_2m,x_1m)\ }M^2
	\xrightarrow{\ (u,v)\mapsto x_1u+x_2v\ }M
	\longrightarrow0.
\]
因此第一同调测量关系 \(x_1u+x_2v=0\) 是否都来自显式的交换关系 \((-x_2m,x_1m)\)。定义和这些低次形式可对照\cite[第16节，第127—128页]{Matsumura1989}。

\begin{proposition}\label{g5:prop:koszul-null-homotopy}
	设 \(A\) 为交换含幺环，\(\mathbf x=(x_1,\ldots,x_r)\)，\(M\) 为 \(A\) 模。对每个 \(j\)，乘以 \(x_j\) 在 \(K_\bullet(\mathbf x;M)\) 上与零链同伦。因而
	\[
		(x_1,\ldots,x_r)H_i(\mathbf x;M)=0
		\qquad\text{对所有 }i.
	\]
\end{proposition}
\begin{proof}
	由微分的定义，在纯楔积上按因子来自前段或后段分别求和，可得分次乘法公式
	\[
		\mathrm d(\omega\wedge\eta)
		=\mathrm d\omega\wedge\eta
		+(-1)^{\deg\omega}\omega\wedge\mathrm d\eta.
	\]
	取 \(\omega=e_j\)，并张量 \(M\)，便得到
	\[
		\mathrm d(e_j\wedge v)=x_jv-e_j\wedge\mathrm d v.
	\]
	所以次数增加一的映射 \(h_j(v)=e_j\wedge v\) 满足
	\(\mathrm d h_j+h_j\mathrm d=x_j\operatorname{id}\)，正是所求链同伦。对循环元 \(v\)，有 \(x_jv=\mathrm d h_j(v)\)，故乘以 \(x_j\) 在同调上为零。
\end{proof}

若某个 \(x_j\) 为单位，将上述同伦除以 \(x_j\)，便知恒等映射与零同伦，全部同调连同 \(H_0\) 都为零。但此时末端商为零，序列不正则。因此同调判别必须同时注意末端非零条件。

\subsection{正则序列给出消解}
\label{g5:subsec:0802:03}

每增加一个元素，Koszul 复形都会由较短复形的两份副本组成。它们之间的乘法映射给出一个长正合列，既能证明正则序列的充分性，也将用于下一小节的反向判别。

\begin{lemma}\label{g5:lem:koszul-mapping-cone}
	设 \(A\) 为交换含幺环，\(M\) 为 \(A\) 模，\(r\geq1\)，并记
	\(\mathbf x'=(x_1,\ldots,x_{r-1})\)、\(x=x_r\)、\(C_\bullet=K_\bullet(\mathbf x';M)\)。则存在自然的同调长正合列
	\begin{equation}\label{g5:eq:koszul-long-exact}
		\begin{aligned}
		\cdots\longrightarrow H_i(C)&\xrightarrow{x}H_i(C)
		\longrightarrow H_i(\mathbf x;M)\\
		&\longrightarrow H_{i-1}(C)\xrightarrow{x}H_{i-1}(C)
		\longrightarrow\cdots.
		\end{aligned}
	\end{equation}
\end{lemma}
\begin{proof}
	把最后一个基向量放在楔积的最前端，每个次数 \(i\) 的元素唯一写成 \(u+e_r\wedge v\)，其中 \(u\in C_i\)、\(v\in C_{i-1}\)。由刚才的分次乘法公式，
	\[
		K_i(\mathbf x;M)=C_i\oplus C_{i-1},
		\qquad
		\mathrm d(u,v)=(\mathrm d u+xv,-\mathrm d v).
	\]
	这就是链映射 \(x:C_\bullet\to C_\bullet\) 的\term[yingshezhui]{映射锥}[mapping cone]。明确地说，令 \(C[1]_i=C_{i-1}\)，其微分取 \(-\mathrm d_C\)，则包含和投影给出逐项分裂的短正合复形列
	\[
		0\longrightarrow C_\bullet
		\xrightarrow{u\mapsto(u,0)}K_\bullet(\mathbf x;M)
		\xrightarrow{(u,v)\mapsto v}C[1]_\bullet
		\longrightarrow0.
	\]
	投影是链映射，正是因为商复形的微分为 \(-\mathrm d_C\)。

	由\cref{g5:prop:complex-long-exact}得到长正合同调列，且 \(H_i(C[1])=H_{i-1}(C)\)。还需辨认连接映射。对循环元 \(v\in C_{i-1}\)，在中间项中选提升 \((0,v)\)，其微分为
	\[
		\mathrm d(0,v)=(xv,0).
	\]
	按照连接映射的定义，其像就是 \([xv]\)。因此连接映射恰为乘以 \(x\)，将各项写开便得到\eqref{g5:eq:koszul-long-exact}。
\end{proof}

\begin{theorem}\label{g5:thm:regular-koszul-acyclic}
	设 \(A\) 为交换含幺环，\(M\) 为 \(A\) 模，\(\mathbf x=(x_1,\ldots,x_r)\) 为 \(M\) 正则序列。则
	\[
		H_i(\mathbf x;M)=0\quad(i>0),\qquad
		H_0(\mathbf x;M)=M/(\mathbf x)M.
	\]
	若 \(M=A\)，则增广复形
	\[
		0\longrightarrow K_r(\mathbf x;A)\longrightarrow\cdots
		\longrightarrow K_0(\mathbf x;A)
		\longrightarrow A/(\mathbf x)\longrightarrow0
	\]
	是 \(A/(\mathbf x)\) 的有限自由消解。
\end{theorem}
\begin{proof}
	对序列长度 \(r\) 归纳。\(r=0\) 时复形只在次数零有 \(M\)，结论成立。若 \(r\geq1\)，前段 \(\mathbf x'\) 为 \(M\) 正则序列，由归纳假设，其复形 \(C_\bullet\) 的所有正次同调为零，且
	\[
		H_0(C)=M/(x_1,\ldots,x_{r-1})M.
	\]
	对 \(i\geq2\)，\eqref{g5:eq:koszul-long-exact}中位于 \(H_i(\mathbf x;M)\) 两侧的 \(H_i(C)\)、\(H_{i-1}(C)\) 都为零，因此中间项为零。对 \(i=1\)，它给出
	\[
		0\longrightarrow H_1(\mathbf x;M)
		\longrightarrow H_0(C)\xrightarrow{x_r}H_0(C).
	\]
	正则性保证最后一张乘法映射单射，故第一同调也为零。零次同调已经直接算出。若 \(M=A\)，各项 \(K_i\) 都是有限自由模，增广的核由 \(\mathrm d_1\) 的像给出，于是得到自由消解。
\end{proof}

这个结论不要求 \(A\) 局部，也不要求 \(M\) 有限生成；这里使用的输入就是逐次非零因子条件。相关长正合列与充分性证明见\cite[第16节，定理16.4—16.5，第127—129页]{Matsumura1989}。

\begin{example}\label{g5:exm:koszul-plane-origin}
	设 \(A=k[x,y]\)。序列 \(x,y\) 正则，故原点的坐标模 \(k=A/(x,y)\) 有自由消解
	\[
		0\longrightarrow A
		\xrightarrow{\binom{-y}{x}}A^2
		\xrightarrow{\ (x\ \ y)\ }A
		\longrightarrow k\longrightarrow0.
	\]
	中间正合性可以直接复算：若 \(xu+yv=0\)，模去 \(x\) 后得到 \(y\bar v=0\) 于 \(k[y]\)，故 \(v=xw\)；代回并约去非零元素 \(x\)，得到 \(u=-yw\)。因此所有关系都由 \((-y,x)\) 生成，正是二次 Koszul 项的像。
\end{example}

\begin{corollary}\label{g5:cor:koszul-computes-tor}
	设 \(A\) 为交换含幺环，\(\mathbf x\) 为 \(A\) 正则序列，\(N\) 为任意 \(A\) 模。则对所有 \(i\geq0\)，有自然同构
	\[
		\operatorname{Tor}_i^A(A/(\mathbf x),N)
		\cong H_i(\mathbf x;N).
	\]
\end{corollary}
\begin{proof}
	把上个定理给出的自由消解张量 \(N\)，得到的复形正是 \(K_\bullet(\mathbf x;N)\)。按\cref{g5:def:ext-tor}中 Tor 的定义，其同调就是左侧；消解选择无关性保证该识别计算的是同一个 Tor。
\end{proof}

必须保留推论中“\(A\) 正则”的条件。若只知道 \(\mathbf x\) 为 \(M\) 正则序列，\(K_\bullet(\mathbf x;M)\) 虽然在正次数正合，各项却只是 \(M\) 的有限直和，未必是自由 \(A\) 模；而 \(K_\bullet(\mathbf x;A)\) 甚至未必正合。例如
\[
	A=k[u,v]/(uv),\qquad M=A/(u)\cong k[v].
\]
元素 \(v\) 是 \(M\) 上的正则序列，末端商为 \(k\)；但 \(v\) 在 \(A\) 上零化非零元素 \(u\)，所以 \(H_1(v;A)\ne0\)。不能把这两种不同系数的复形混作同一个自由消解。

\subsection{局部有限模上的反向正合判据}
\label{g5:subsec:0802:04}

反过来，要从同调消失恢复逐次单射，需要防止较短复形的非零同调被最后一个元素全部“移走”。局部环中的极大理想条件与有限性，使 Nakayama 引理恰好能够排除这种情况。

\begin{theorem}\label{g5:thm:koszul-regularity-criterion}
	设 \((A,\mathfrak m)\) 为交换含幺 Noether 局部环，\(M\ne0\) 为有限生成 \(A\) 模，\(\mathbf x=(x_1,\ldots,x_r)\) 的每个元素都属于 \(\mathfrak m\)。则以下条件等价：
	\begin{equivalences}
		\item \(\mathbf x\) 为 \(M\) 正则序列；
		\item \(H_i(\mathbf x;M)=0\) 对每个 \(i>0\) 成立；
		\item \(H_1(\mathbf x;M)=0\)。
	\end{equivalences}
\end{theorem}
\begin{proof}
	第一条件蕴含第二条件，已由\cref{g5:thm:regular-koszul-acyclic}证明；第二条件显然蕴含第三条件。现假设第三条件成立，对 \(r\) 归纳证明正则性。\(r=0\) 时为空序列，\(M\ne0\) 保证结论。

	设 \(r\geq1\)，取 \(C_\bullet=K_\bullet(x_1,\ldots,x_{r-1};M)\)。\eqref{g5:eq:koszul-long-exact}中的一段为
	\[
		H_1(C)\xrightarrow{x_r}H_1(C)
		\longrightarrow H_1(\mathbf x;M)=0.
	\]
	于是 \(H_1(C)=x_rH_1(C)\)。每个 \(C_i\) 都是有限模 \(M\) 的有限直和；因 \(A\) Noether，其循环子模有限生成，同调作为该子模的商也有限生成。又 \(x_r\in\mathfrak m\)，故
	\[
		H_1(C)=x_rH_1(C)\subseteq\mathfrak mH_1(C)\subseteq H_1(C).
	\]
	Nakayama 引理给出 \(H_1(C)=0\)。由归纳假设，前段 \(x_1,\ldots,x_{r-1}\) 为 \(M\) 正则序列。

	同一长正合列继续给出
	\[
		0=H_1(\mathbf x;M)\longrightarrow
		H_0(C)\xrightarrow{x_r}H_0(C).
	\]
	因此 \(x_r\) 在 \(H_0(C)=M/(x_1,\ldots,x_{r-1})M\) 上单射。最后，\cref{g5:prop:regular-sequence-terminal-quotient}保证 \(M/(\mathbf x)M\ne0\)，所以完整序列满足定义。
\end{proof}

证明中，Noether 条件保证所需的同调模有限生成，极大理想条件使 Nakayama 引理可用，\(M\ne0\) 则保证末端商仍非零。本节采用的局部 Noether 有限模判据，可对照\cite[第16节，定理16.5，第128—129页]{Matsumura1989}。这些条件也解释了为何正次同调的消失本身，不能在任意环与任意模上直接当作正则序列的定义。

\subsection{交换次序和基变换的限制}
\label{g5:subsec:0802:05}

外代数构造允许对记录方程的自由模换基。复形在这种变化下始终保持同构，而把同构转回“逐次非零因子”的结论，则需要刚才的反向判据。

\begin{proposition}\label{g5:prop:koszul-change-basis}
	设 \(A\) 为交换含幺环，\(M\) 为 \(A\) 模，\(\mathbf x=(x_1,\ldots,x_r)\)，\(U=(u_{ij})\in\operatorname{GL}_r(A)\)。令
	\[
		y_j=\sum_{i=1}^r x_i u_{ij}\qquad(1\leq j\leq r).
	\]
	则存在链同构 \(K_\bullet(\mathbf y;M)\cong K_\bullet(\mathbf x;M)\)。特别地，任意置换元素次序都给出同构的 Koszul 复形。
\end{proposition}
\begin{proof}
	在一次自由模上，把 \(\mathbf y\) 对应的基向量 \(e'_j\) 送到 \(\sum_i u_{ij}e_i\)，再对各次外幂取诱导映射，并张量 \(\operatorname{id}_M\)。矩阵 \(U\) 可逆，故各次映射均可逆。在次数一有
	\[
		\mathrm d_{\mathbf x}\left(\sum_i u_{ij}e_i\right)
		=\sum_i u_{ij}x_i=y_j.
	\]
	利用微分的分次乘法公式，在基楔积上逐项应用这个等式，就知各次映射与微分交换。置换矩阵是可逆矩阵，故包含所述特殊情形。
\end{proof}

\begin{corollary}\label{g5:cor:regular-sequence-permutation}
	设 \((A,\mathfrak m)\) 为 Noether 局部环，\(M\ne0\) 为有限生成 \(A\) 模，\(x_1,\ldots,x_r\in\mathfrak m\) 为 \(M\) 正则序列。则任何置换后的序列仍为 \(M\) 正则序列。更一般地，按上个命题用任意可逆矩阵 \(U\) 改换这组元素，所得 \(y_1,\ldots,y_r\) 仍为 \(M\) 正则序列。
\end{corollary}
\begin{proof}
	原 Koszul 复形的第一同调为零，换基后的第一同调因链同构而仍为零。各 \(y_j\) 都属于 \(\mathfrak m\)，故\cref{g5:thm:koszul-regularity-criterion}给出所需正则性。
\end{proof}

\begin{example}\label{g5:exm:nonlocal-regular-order}
	设 \(k\) 为域，\(A=k\times k[x,y]\)，取
	\[
		a=(1,x),\qquad b=(0,y).
	\]
	乘以 \(a\) 在两个分量上都单射，且
	\[
		A/aA\cong k[y].
	\]
	在这个商环上，\(b\) 的作用就是乘以 \(y\)，仍单射；末端商 \(A/(a,b)\cong k\ne0\)。所以 \(a,b\) 是 \(A\) 正则序列。但 \(b\) 已经零化非零元素 \((1,0)\)，故 \(b,a\) 不是正则序列。两个次序的 Koszul 复形仍然同构，并且正次同调都消失；失效的是非局部情形下从同调消失恢复逐次单射的推理。
\end{example}

可逆换基也不能替换成任意增添生成元。以 \(A=k[x]\) 为例，\(x\) 是正则序列，而 \(x,x\) 生成同一个理想却不正则：第二个 \(x\) 在 \(A/(x)\) 上为零。其第一 Koszul 同调也非零，因为
\[
	\ker\bigl((x\ \ x):A^2\to A\bigr)=A(1,-1),
	\qquad
	\operatorname{im}\mathrm d_2=xA(1,-1),
\]
故 \(H_1(x,x;A)\cong A/(x)\)。因此 Koszul 同调记录所选元素组的关系，不只记录它们生成的理想。

\begin{proposition}\label{g5:prop:koszul-base-change}
	设 \(A\to B\) 为交换含幺环同态，\(M\) 为 \(A\) 模，\(\mathbf x=(x_1,\ldots,x_r)\)，\(\mathbf x_B\) 表示这些元素在 \(B\) 中的像。则有自然链同构
	\[
		K_\bullet(\mathbf x;M)\otimes_A B
		\cong K_\bullet(\mathbf x_B;M\otimes_A B).
	\]
	若 \(B\) 对 \(A\) 平坦，则还存在自然同构
	\[
		H_i(\mathbf x;M)\otimes_A B
		\cong H_i(\mathbf x_B;M\otimes_A B)
		\qquad\text{对所有 }i.
	\]
\end{proposition}
\begin{proof}
	第一项只需按有限自由外幂的基作标量扩张：两边的基楔积一一对应，微分矩阵的条目都是带符号的 \(x_j\)，变换系数后正好得到右侧微分，所以不需要平坦性。

	若 \(B\) 平坦，对原复形记循环模 \(Z_i\)、边界模 \(D_i\)。把短正合列
	\[
		0\longrightarrow Z_i\longrightarrow K_i\longrightarrow D_{i-1}\longrightarrow0,
		\qquad
		0\longrightarrow D_i\longrightarrow Z_i\longrightarrow H_i\longrightarrow0
	\]
	张量 \(B\) 后仍短正合。再用 \(D_{i-1}\hookrightarrow K_{i-1}\) 张量后的单射，可知新复形的循环模为 \(Z_i\otimes_A B\)，边界模为 \(D_i\otimes_A B\)。两者取商即得同调同构。
\end{proof}

\begin{proposition}\label{g5:prop:regular-sequence-flat-base-change}
	设 \(A\to B\) 为平坦的交换含幺环同态，\(M\) 为 \(A\) 模，\(\mathbf x\) 为 \(M\) 正则序列。若
	\[
		\bigl(M/(\mathbf x)M\bigr)\otimes_A B\ne0,
	\]
	则 \(\mathbf x_B\) 为 \(M\otimes_A B\) 正则序列。若 \(A\to B\) 忠实平坦，则对任意有限元素组 \(\mathbf x\)，原序列为 \(M\) 正则序列，当且仅当其像为 \(M\otimes_A B\) 正则序列。
\end{proposition}
\begin{proof}
	对 \(M_j=M/(x_1,\ldots,x_j)M\)，张量积右正合性给出
	\[
		M_j\otimes_A B
		\cong (M\otimes_A B)/(x_1,\ldots,x_j)(M\otimes_A B).
	\]
	逐次乘法单射在平坦张量后仍单射，末端非零由假设保证，故正则性保持。

	若同态忠实平坦，非零末端商也不会在张量后消失，所以正方向不再需要附加假设。反过来，基变换后各次乘法单射，由\cref{g5:prop:faithful-flat-reflects-exactness}反映为原来的单射；基变换后的末端商非零，也迫使原末端商非零。于是原序列正则。
\end{proof}

两个限制都能在单个元素上出现。对 \(A=k[x]\) 上的正则元素 \(x\)，取非平坦商 \(B=A/(x)=k\)，则基变换后的元素为零，乘法不再单射；原来的第一 Koszul 同调为零，新的第一同调却是 \(k\)。若改取平坦局部化 \(B=A[x^{-1}]\)，乘法仍单射，但 \(x\) 成为单位，末端商消失，因此也不再是这里定义的正则序列。平坦性控制单射，忠实性或明确的末端非零条件控制剩下的商模，这两项作用不可互相替代。
% 纲要条目排为编号小节标题，并进入目录与书签；说明段仍保留于纲要注释中。

% !TeX root = ../../Book5.tex
% 纲要标题：### 8.3 深度的不同刻画
% 纲要位置：几何学计划纲要.md，第 5319 行。
% 正文状态：初稿；尚未正式审核。

% 纲要内容开始
% 
% * 极大正则序列长度
% * Ext 检测与等价证明
% * 深度不超过维数
% * 深度引理
% * Serre Sᵣ 条件的局部表述
% 
% 纲要内容结束

\section{深度的不同刻画}
\label{g5:sec:0803}

在一个局部环上，选择正则元相当于用一个方程取商，同时要求乘法保持单射。连续这样做，能够进行多少次？支集的维数给出上界，但嵌入素理想可能使过程提前停止。我们要把这个停止位置变成模的不变量，并证明它与选择哪些正则元无关。

\ProseParagraphBreak

本节先证明每一步都使支集维数下降，保证正则序列能延长到某个终点。随后用 \(\operatorname{Ext}^i_A(k,M)\) 检测这个终点：它们起初连续消失的项数，恰好等于正则序列能够达到的长度。长正合列于是把深度转化为短正合列中的可比较数据，逐素局部化又给出 Serre 条件。正则序列与 Ext 刻画可对照\cite[第16节，定理16.6—16.7，第129—131页]{Matsumura1989}。

\subsection{极大正则序列长度}
\label{g5:subsec:0803:01}

\begin{definition}\label{g5:def:module-dimension}
	设 \(A\) 为 Noether 环，\(M\) 为有限生成 \(A\) 模。\term[moweishu]{模的维数}[dimension of a module]定义为其支集的维数：
	\[
		\dim_A M=\dim\operatorname{Supp}_A(M)
		=\dim\bigl(A/\operatorname{Ann}_A(M)\bigr).
	\]
	这里支集维数按支集中严格素链的长度定义；第二个等号由\cref{g5:prop:support-annihilator}及商环的素理想对应得到。约定 \(\dim_A0=-\infty\)。
\end{definition}

若 \((A,\mathfrak m)\) 为 Noether 局部环，\(M\ne0\) 有限生成，则 \(\operatorname{Ann}_A(M)\subseteq\mathfrak m\)，且 \(\dim_A M\) 是有限的非负整数。事实上，\(\mathfrak m\) 有有限个生成元，高度定理 \cref{g5:thm:height-theorem}便使 \(\dim A=\operatorname{ht}_A\mathfrak m\) 有限，商环的维数不会更大。

\begin{lemma}\label{g5:lem:regular-element-support-drop}
	设 \((A,\mathfrak m)\) 为 Noether 局部环，\(M\ne0\) 为有限生成 \(A\) 模，且 \(x\in\mathfrak m\) 是 \(M\) 正则元。则 \(M/xM\ne0\)，并且
	\[
		\operatorname{Supp}_A(M/xM)
		=\operatorname{Supp}_A(M)\cap V(x),
		\qquad
		\dim_A(M/xM)\leq\dim_A M-1.
	\]
\end{lemma}
\begin{proof}
	因为 \(x\in\mathfrak m\)，Nakayama 引理 \cref{g5:thm:nakayama-local}说明 \(M/xM\ne0\)。对任意素理想 \(\mathfrak p\)，局部化与取商相容，给出
	\[
		(M/xM)_{\mathfrak p}
		=M_{\mathfrak p}/xM_{\mathfrak p}.
	\]
	若 \(M_{\mathfrak p}=0\)，或者 \(x\notin\mathfrak p\) 而在 \(A_{\mathfrak p}\) 中成为单位，这个商为零。若 \(M_{\mathfrak p}\ne0\) 且 \(x\in\mathfrak p\)，则再次在局部环 \(A_{\mathfrak p}\) 上用 Nakayama 引理，得到这个商非零。这证明支集等式；这一等式本身并不需要 \(x\) 正则。

	由\cref{g5:cor:minimal-support-associated}，支集的极小素理想都是 \(M\) 的关联素理想。\cref{g5:thm:zero-divisors-associated}又说明，正则元 \(x\) 不属于其中任何一个。现在取商模支集中的任意严格素链
	\[
		\mathfrak q_0\subsetneq\mathfrak q_1
		\subsetneq\cdots\subsetneq\mathfrak q_l.
	\]
	在 \(\mathfrak q_0\) 下方选择 \(\operatorname{Supp}_A(M)\) 的一个极小元 \(\mathfrak p\)。由于 \(x\in\mathfrak q_0\) 而 \(x\notin\mathfrak p\)，有 \(\mathfrak p\subsetneq\mathfrak q_0\)。将它接在链的最前面，就得到 \(\operatorname{Supp}_A(M)\) 中长度为 \(l+1\) 的链。因此 \(l+1\leq\dim_A M\)，对所有这样的链取上界即得结论。
\end{proof}

回顾\cref{g5:def:module-regular-sequence}，位于 \(\mathfrak m\) 中的 \(M\) 正则序列 \(x_1,\ldots,x_t\) 要求每个 \(x_i\) 在
\(M/(x_1,\ldots,x_{i-1})M\) 上的乘法为单射，且最终商非零。对这里的局部有限模，各中间商的非零性也由 Nakayama 引理保证。逐次应用上面的引理，序列长度必满足 \(t\leq\dim_A M\)。所以从空序列出发，只要还能追加一个正则元就继续，至多经过有限步便会停止。

\begin{definition}\label{g5:def:module-depth}
	设 \((A,\mathfrak m)\) 为 Noether 局部环，\(M\ne0\) 为有限生成 \(A\) 模。若 \(\mathfrak m\) 中的 \(M\) 正则序列不能再在末尾追加一个 \(\mathfrak m\) 中的元素而仍保持正则，就称它为极大正则序列。所有 \(\mathfrak m\) 中的 \(M\) 正则序列长度的最大值，称为 \(M\) 的\term[shendu]{深度}[depth]，记为 \(\operatorname{depth}_A M\)。空序列的长度为零，并另约定
	\[
		\operatorname{depth}_A0=\infty.
	\]
\end{definition}
\symidx{\(\operatorname{depth}_A M\)：有限模的深度}

最大长度存在，来自刚才证明的整数上界。“极大”目前只表示不能继续追加；下一小节将证明，任意极大正则序列都达到这个最大长度。首先确定过程一步也不能开始的情形。以下 \(k=A/\mathfrak m\) 表示剩余域。

\begin{lemma}\label{g5:lem:depth-zero}
	设 \((A,\mathfrak m,k)\) 为 Noether 局部环，\(M\ne0\) 为有限生成 \(A\) 模。以下条件等价：
	\begin{equivalences}
		\item \(\operatorname{depth}_A M=0\)；
		\item \(\mathfrak m\in\operatorname{Ass}_A(M)\)；
		\item \(\operatorname{Hom}_A(k,M)\ne0\)。
	\end{equivalences}
\end{lemma}
\begin{proof}
	第一个条件表示 \(\mathfrak m\) 中每个元素都是 \(M\) 零因子。由\cref{g5:thm:prime-filtration,g5:thm:zero-divisors-associated}，关联素理想有限，且这些零因子恰为它们的并。有限素理想避开引理 \cref{g5:lem:finite-prime-avoidance}表明，\(\mathfrak m\) 包含在某个关联素理想中。局部环的所有真素理想又都包含在 \(\mathfrak m\) 中，故该关联素理想只能是 \(\mathfrak m\)。反过来，若 \(\mathfrak m\) 本身关联，它的每个元素都是零因子，所以没有长度一的正则序列。

	同态 \(k\to M\) 由 \(1\) 的像 \(m\) 唯一确定，而这个像必须满足 \(\mathfrak m m=0\)。非零同态的像 \(m\ne0\) 的零化子是真理想，包含极大理想 \(\mathfrak m\)，故恰等于 \(\mathfrak m\)。这与第二个条件等价。
\end{proof}

\subsection{Ext 检测与等价证明}
\label{g5:subsec:0803:02}

固定局部环 \((A,\mathfrak m,k)\)，暂记
\[
	E^i(M)=\operatorname{Ext}^i_A(k,M)\qquad(i\geq0).
\]
\cref{g5:prop:ext-annihilator}说明，\(\mathfrak m\) 零化每个 \(E^i(M)\)：第一变量 \(k\) 上的标量作用，与第二变量上的标量作用诱导同一作用。因此对于 \(x\in\mathfrak m\)，由第二变量乘以 \(x\) 诱导的全部 \(E^i(M)\to E^i(M)\) 都为零。这正好把正则元对应的长正合列拆成短正合列。

\begin{lemma}\label{g5:lem:regular-element-ext-sequence}
	设 \((A,\mathfrak m,k)\) 为 Noether 局部环，\(M\ne0\) 为有限生成 \(A\) 模，\(x\in\mathfrak m\) 为 \(M\) 正则元，并令 \(\overline M=M/xM\)、\(E^i(T)=\operatorname{Ext}^i_A(k,T)\)。则 \(E^0(M)=0\)，且对每个 \(i\geq0\) 有短正合列
	\begin{equation}\label{g5:eq:regular-ext-short-exact}
		0\longrightarrow E^i(M)\longrightarrow E^i(\overline M)
		\longrightarrow E^{i+1}(M)\longrightarrow0.
	\end{equation}
\end{lemma}
\begin{proof}
	正则性给出短正合列
	\[
		0\longrightarrow M\xrightarrow{\,x\,}M
		\longrightarrow\overline M\longrightarrow0.
	\]
	由\cref{g5:prop:ext-tor-long-exact}，其 Ext 长正合列包含
	\[
		E^i(M)\xrightarrow{\,x\,}E^i(M)
		\longrightarrow E^i(\overline M)
		\longrightarrow E^{i+1}(M)\xrightarrow{\,x\,}E^{i+1}(M).
	\]
	两张标为 \(x\) 的映射均为零，便得到\eqref{g5:eq:regular-ext-short-exact}。在长正合列最左端，\(E^0(M)\xrightarrow{x}E^0(M)\) 同时为单射和零映射，所以 \(E^0(M)=0\)。
\end{proof}

\begin{theorem}\label{g5:thm:depth-ext}
	设 \((A,\mathfrak m,k)\) 为 Noether 局部环，\(M\ne0\) 为有限生成 \(A\) 模。则
	\[
		\operatorname{depth}_A M
		=\min\{i\geq0:\operatorname{Ext}^i_A(k,M)\ne0\}.
	\]
	右侧集合非空，而且 \(\mathfrak m\) 中任意极大 \(M\) 正则序列的长度都等于这个整数。对零模，采用 \(\inf\varnothing=\infty\)，同一刻画仍成立。
\end{theorem}
\begin{proof}
	先观察\eqref{g5:eq:regular-ext-short-exact}如何改变首个非零次数。若 \(x\) 正则，且 \(E^i(M/xM)\) 的首个非零次数为 \(r\)，则 \(E^0(M)=0\)。对 \(i<r\)，短正合列中间为零，使两端都为零，所以
	\[
		E^0(M)=\cdots=E^r(M)=0.
	\]
	在 \(i=r\) 时，同一短正合列给出
	\[
		E^{r+1}(M)\cong E^r(M/xM)\ne0.
	\]
	因此首个非零次数向上增加一；\(r=0\) 时也由最左端消失及 \(i=0\) 的短正合列得到这一结论。

	现在取任意极大 \(M\) 正则序列 \(x_1,\ldots,x_t\)，并记
	\[
		M_j=M/(x_1,\ldots,x_j)M\qquad(0\leq j\leq t).
	\]
	这些模均非零且有限生成。极大性说明 \(M_t\) 没有 \(\mathfrak m\) 中的正则元，由\cref{g5:lem:depth-zero}，有 \(E^0(M_t)\ne0\)。从 \(j=t\) 向前逐次使用刚才的观察，得到 \(E^i(M)\) 的首个非零次数恰为 \(t\)。

	右侧由 \(M\) 本身决定，所以每条极大序列的长度都相同。任意正则序列都能在有限步后延长到极大序列，因而这个共同长度就是定义中的最大长度。对于零模，全部 Ext 群为零，所作无穷大约定与之相容。
\end{proof}

这个证明同时解释了 Ext 为何能够检测正则序列：每次对正则元取商，都会把第一次非零的 Ext 向低一次移动，而终点由 \(\operatorname{Hom}_A(k,-)\) 检出。这与\cite[第16节，定理16.7，第130页]{Matsumura1989}的深度刻画一致。

\begin{corollary}\label{g5:cor:regular-depth-drop}
	设 \((A,\mathfrak m)\) 为 Noether 局部环，\(M\ne0\) 为有限生成 \(A\) 模，且 \(x_1,\ldots,x_t\in\mathfrak m\) 为 \(M\) 正则序列。则
	\[
		\operatorname{depth}_A\bigl(M/(x_1,\ldots,x_t)M\bigr)
		=\operatorname{depth}_A M-t.
	\]
	左侧也等于将该商模看作 \(A/(x_1,\ldots,x_t)\) 模时的深度。
\end{corollary}
\begin{proof}
	对 \(t=1\)，将 \(M/x_1M\) 的极大正则序列接在 \(x_1\) 之后，就得到 \(M\) 的极大正则序列；由定理，它们的长度相差一。逐个取商即得一般结论。若理想 \(H\) 零化一个模，\(A\) 与 \(A/H\) 中对应元素对该模及逐次商模的乘法作用完全相同，故正则序列在取像与提升下对应。这给出最后的环变换说明。
\end{proof}

\subsection{深度不超过维数}
\label{g5:subsec:0803:03}

现在可以把正则序列的长度与实际削去的支集维数放在一起比较。极小支集素理想的避开已经给出一个方向；主理想的维数下降不等式补出另一个方向。

\begin{theorem}\label{g5:thm:depth-dimension}
	设 \((A,\mathfrak m)\) 为 Noether 局部环，\(M\ne0\) 为有限生成 \(A\) 模。则
	\[
		0\leq\operatorname{depth}_A M\leq\dim_A M\leq\dim A.
	\]
	若 \(x_1,\ldots,x_t\in\mathfrak m\) 为 \(M\) 正则序列，则
	\[
		\dim_A\bigl(M/(x_1,\ldots,x_t)M\bigr)=\dim_A M-t.
	\]
\end{theorem}
\begin{proof}
	第一串不等式来自正则序列长度的支集维数上界，以及 \(\operatorname{Supp}_A(M)\) 是 \(\operatorname{Spec}A\) 的闭子集。

	为证明精确下降公式，先取一个 \(M\) 正则元 \(x\in\mathfrak m\)，并记 \(H=\operatorname{Ann}_A(M)\)。由\cref{g5:lem:regular-element-support-drop}中的支集等式及 \(\operatorname{Supp}_A(M)=V(H)\)，有
	\[
		\dim_A(M/xM)=\dim A/(H+(x)).
	\]
	环 \(A/H\) 是 Noether 局部环，\(x\) 的像属于其极大理想。对它使用\cref{g5:lem:principal-dimension-drop}，得到
	\[
		\dim_A M=\dim A/H
		\leq\dim A/(H+(x))+1
		=\dim_A(M/xM)+1.
	\]
	再结合\cref{g5:lem:regular-element-support-drop}的相反不等式，得到 \(\dim_A(M/xM)=\dim_A M-1\)。对各非零中间商重复应用，就得到长度 \(t\) 的公式。
\end{proof}

维数描述支集还能容纳多长的素链，深度则要求每一步取商前的乘法都保持单射。两者相等时，极大正则序列恰好把支集降到零维；若不相等，正则序列会在仍有正维支集时停止。相关维数比较也可对照\cite[第16节，习题16.1，第132页；第17节，定理17.2，第134页]{Matsumura1989}。

\begin{example}\label{g5:exm:depth-support-dimension}
	设 \(k\) 为域，\(A=k[x,y]_{(x,y)}\)。\(x,y\) 是 \(A\) 正则序列，商为 \(k\)，而 \(\dim A=2\)，故 \(\operatorname{depth}A=2\)。对模 \(M=A/(x)\)，有
	\[
		M\cong k[y]_{(y)},\qquad \dim_A M=1.
	\]
	\(y\) 在 \(M\) 上正则，且 \(M/yM\cong k\)，所以 \(\operatorname{depth}_A M=1\)。这里与深度相比较的维数是模的支集维数，而非环境环的维数 \(2\)。

	再取 \(N=A/(x^2,xy)\)。其支集仍是 \(V(x)\)，故 \(\dim_A N=1\)。但非零元素 \(\overline x\) 被 \((x,y)\) 零化，于是 \(\operatorname{Hom}_A(k,N)\ne0\)，得到 \(\operatorname{depth}_A N=0\)。\(\overline x\) 在局部化后仍非零：否则存在多项式 \(s\in k[x,y]\setminus(x,y)\)，使 \(sx\in(x^2,xy)\)；在多项式环中约去 \(x\) 会迫使 \(s\in(x,y)\)，矛盾。这个模的关联素理想也已经在\cref{g5:exm:associated-embedded-point}中计算；嵌入于原点的部分使正则取商无法开始。
\end{example}

零模的约定与上面的不等式分开使用：\(\operatorname{depth}0=\infty\)、\(\dim0=-\infty\)，所以深度不超过维数这一断言只针对非零有限模。

\subsection{深度引理}
\label{g5:subsec:0803:04}

若一个模由子模和商模拼成，支集只是两者支集的并，而深度还受连接映射影响。Ext 长正合列给出三个方向的估计，并在深度不相等时给出精确等式。

\begin{lemma}[深度引理]\label{g5:lem:depth-lemma}
	设 \((A,\mathfrak m)\) 为 Noether 局部环，且
	\[
		0\longrightarrow L\longrightarrow M\longrightarrow N\longrightarrow0
	\]
	是有限生成 \(A\) 模的短正合列。则
	\begin{align*}
		\operatorname{depth}M
		&\geq\min\{\operatorname{depth}L,\operatorname{depth}N\},\\
		\operatorname{depth}L
		&\geq\min\{\operatorname{depth}M,\operatorname{depth}N+1\},\\
		\operatorname{depth}N
		&\geq\min\{\operatorname{depth}L-1,\operatorname{depth}M\}.
	\end{align*}
	其中深度都相对于 \(A\)，并采用 \(\infty+1=\infty-1=\infty\)。
\end{lemma}
\begin{proof}
	记 \(a=\operatorname{depth}L\)、\(b=\operatorname{depth}M\)、\(c=\operatorname{depth}N\)，仍写 \(E^i(T)=\operatorname{Ext}^i_A(k,T)\)，并令 \(E^i(T)=0\) 当 \(i<0\)。对短正合列使用\cref{g5:prop:ext-tor-long-exact}。

	若 \(i<\min(a,c)\)，长正合列的片段
	\[
		E^i(L)\longrightarrow E^i(M)\longrightarrow E^i(N)
	\]
	两端为零，所以 \(E^i(M)=0\)。由\cref{g5:thm:depth-ext}，\(b\geq\min(a,c)\)。若 \(i<\min(b,c+1)\)，则
	\[
		E^{i-1}(N)\longrightarrow E^i(L)\longrightarrow E^i(M)
	\]
	两端为零，得到 \(a\geq\min(b,c+1)\)。最后，若 \(i<\min(a-1,b)\)，则
	\[
		E^i(M)\longrightarrow E^i(N)\longrightarrow E^{i+1}(L)
	\]
	两端为零，得到 \(c\geq\min(a-1,b)\)。这些消失论证也涵盖零模及无穷大约定。
\end{proof}

这一结果称为\term[shenduyinli]{深度引理}[depth lemma]。三个不等式分别读取长正合列中相邻三项的消失。若其中两个深度严格不等，其余一项的消失范围会受到两侧估计的共同限制，从而得到下面的等式。

\begin{corollary}\label{g5:cor:depth-lemma-equalities}
	在\cref{g5:lem:depth-lemma}的假设下，记
	\(a=\operatorname{depth}L\)、\(b=\operatorname{depth}M\)、\(c=\operatorname{depth}N\)。则有
	\[
		\begin{array}{c@{\quad\Longrightarrow\quad}c@{\qquad}c@{\quad\Longrightarrow\quad}c}
			a<b & c=a-1 & a>b & c=b,\\
			b<c & a=b & b>c & a=c+1.
		\end{array}
	\]
\end{corollary}
\begin{proof}
	含零模时，短正合列使其余非零项同构，以上结论按无穷大约定直接成立。以下设三项都非零，所以 \(a,b,c\) 都为有限非负整数。

	若 \(a<b\)，由 \(a\geq\min(b,c+1)\) 得 \(c+1\leq a\)，而第三个不等式给出 \(c\geq a-1\)，所以 \(c=a-1\)。若 \(a>b\)，第一个不等式迫使 \(c\leq b\)，而整数不等式 \(a-1\geq b\) 及第三个不等式给出 \(c\geq b\)，故 \(c=b\)。

	若 \(b<c\)，第一个不等式迫使 \(a\leq b\)，第二个又给出 \(a\geq b\)，故 \(a=b\)。若 \(b>c\)，第二个给出 \(a\geq c+1\)，第三个则迫使 \(a-1\leq c\)，故 \(a=c+1\)。
\end{proof}

\begin{example}\label{g5:exm:depth-maximal-ideal}
	取 \(A=k[x,y]_{(x,y)}\)、\(\mathfrak m=(x,y)A\)，其中 \(k\) 为域。短正合列
	\[
		0\longrightarrow\mathfrak m\longrightarrow A\longrightarrow k\longrightarrow0
	\]
	中，\(\operatorname{depth}A=2\)，而 \(\operatorname{depth}k=0\)。由上个推论的 \(b>c\) 情形，\(\operatorname{depth}_A\mathfrak m=1\)。但 \(A\) 是整环且 \(\mathfrak m\ne0\)，所以 \(\operatorname{Ann}_A(\mathfrak m)=0\)，从而 \(\dim_A\mathfrak m=2\)。同一个正合列既改变生成元，也改变了可连续选取正则元的次数。
\end{example}

\subsection{\texorpdfstring{Serre \(S_r\) 条件的局部表述}{Serre Sᵣ 条件的局部表述}}
\label{g5:subsec:0803:05}

在一个指定局部环上计算深度，只检查了一个局部问题。若希望控制模在整个支集上出现的嵌入部分，以及更高维位置上的正则序列，就要在每个素理想处作同样的检测。

\begin{definition}\label{g5:def:serre-module-condition}
	设 \(A\) 为 Noether 环，\(M\) 为有限生成 \(A\) 模，\(r\geq0\) 为整数。若对每个 \(\mathfrak p\in\operatorname{Supp}_A(M)\) 都有
	\[
		\operatorname{depth}_{A_{\mathfrak p}}M_{\mathfrak p}
		\geq\min\{r,\dim_{A_{\mathfrak p}}M_{\mathfrak p}\},
	\]
	则称 \(M\) 满足\term[SerreSr]{Serre \(S_r\) 条件}[Serre's condition \(S_r\)]。对环 \(A\) 本身，就把 \(M\) 取为 \(A\)。零模的支集为空，约定它满足全部 \(S_r\) 条件。
\end{definition}

每个有限模都满足 \(S_0\)，而 \(S_{r+1}\) 蕴含 \(S_r\)。\(S_1\) 要求正维局部支集上的深度至少为一；\(S_2\) 还要求局部支集维数至少为二时，深度至少为二。环的版本把局部支集维数写成 \(\dim A_{\mathfrak p}=\operatorname{ht}_A\mathfrak p\)，见\cite[第23节，第183页]{Matsumura1989}。对一般模，定义中须保留 \(\dim_{A_{\mathfrak p}}M_{\mathfrak p}\)。

\begin{proposition}\label{g5:prop:serre-localization}
	设 \(A\) 为 Noether 环，\(M\) 为有限生成 \(A\) 模，\(r\geq0\)。若 \(M\) 满足 \(S_r\)，则对任意乘法集 \(S\subseteq A\)，\(S^{-1}M\) 作为 \(S^{-1}A\) 模也满足 \(S_r\)。反过来，若每个极大理想 \(\mathfrak m\) 处的 \(M_{\mathfrak m}\) 都作为 \(A_{\mathfrak m}\) 模满足 \(S_r\)，则 \(M\) 满足 \(S_r\)。
\end{proposition}
\begin{proof}
	\(S^{-1}A\) 的支集素点对应于 \(\operatorname{Supp}_A(M)\) 中与 \(S\) 不交的素理想 \(\mathfrak p\)。在该点进一步局部化，得到的局部环和模分别自然同构于 \(A_{\mathfrak p}\)、\(M_{\mathfrak p}\)。所以深度与支集维数都相同，原不等式仍成立。

	反过来，对任意 \(\mathfrak p\in\operatorname{Supp}_A(M)\)，取包含它的极大理想 \(\mathfrak m\)。\(M_{\mathfrak p}\) 是 \(M_{\mathfrak m}\) 在素理想 \(\mathfrak pA_{\mathfrak m}\) 处的局部化。将 \(M_{\mathfrak m}\) 的 \(S_r\) 条件用于这一点，就得到所需不等式。
\end{proof}

\begin{proposition}\label{g5:prop:serre-s1-associated}
	设 \(A\) 为 Noether 环，\(M\) 为有限生成 \(A\) 模。则 \(M\) 满足 \(S_1\)，当且仅当其全部关联素理想都是支集的极小元，即 \(M\) 没有嵌入关联素理想。
\end{proposition}
\begin{proof}
	零模的两边都成立。设 \(M\ne0\)。若 \(\mathfrak p\in\operatorname{Ass}_A(M)\)，由\cref{g5:thm:associated-localization}，\(\mathfrak pA_{\mathfrak p}\) 是 \(M_{\mathfrak p}\) 的关联素理想，因此\cref{g5:lem:depth-zero}给出 \(\operatorname{depth}_{A_{\mathfrak p}}M_{\mathfrak p}=0\)。若满足 \(S_1\)，必须有 \(\dim_{A_{\mathfrak p}}M_{\mathfrak p}=0\)。局部化的素理想对应表明，这恰表示在 \(\mathfrak p\) 下方没有更小的支集素理想，亦即 \(\mathfrak p\) 为极小支集素理想。

	反过来，若每个关联素理想都极小，取 \(\mathfrak p\in\operatorname{Supp}_A(M)\)。若 \(\dim M_{\mathfrak p}=0\)，\(S_1\) 不等式自动成立。若 \(\dim M_{\mathfrak p}>0\)，\(\mathfrak p\) 不是极小支集素理想，故 \(\mathfrak p\notin\operatorname{Ass}_A(M)\)。再次由关联素理想的局部化公式，\(\mathfrak pA_{\mathfrak p}\notin\operatorname{Ass}_{A_{\mathfrak p}}(M_{\mathfrak p})\)，于是其深度至少为一。这正是所需条件。
\end{proof}

\begin{example}\label{g5:exm:serre-support-and-closed-point}
	设 \(A=k[x,y]_{(x,y)}\)，\(M=A/(x)\)，其中 \(k\) 为域。\(M\) 的支集只有 \((x)A\) 和 \((x,y)A\) 两个素点，对应于 \(k[y]_{(y)}\) 的两个素理想。在前一点，局部模的维数和深度均为零；在后一点，两者均为一。因此 \(M\) 满足全部 \(S_r\)。若误用环境局部环的维数 \(\dim A=2\)，就会在极大理想处错误地排除 \(S_2\)。

	再考虑 \(T=A\oplus A/(x)\)。有限直和的 Ext 是各分量 Ext 的直和，这由对第一变量的同一个自由消解取 \(\operatorname{Hom}\) 后逐项分解即可。因此在极大理想处
	\[
		\operatorname{depth}_A T=\min\{2,1\}=1,
		\qquad \dim_A T=2.
	\]
	但在 \(\mathfrak p=(x)A\) 处，有
	\[
		T_{\mathfrak p}\cong A_{\mathfrak p}\oplus\kappa(\mathfrak p).
	\]
	第二个分量使其深度为零，第一个分量使其支集维数为 \(\dim A_{\mathfrak p}=1\)。这里高度为一，是因为 \((x)\) 非零且为素理想，而主理想定理给出其高度至多一。故 \(T\) 不满足 \(S_1\)，尽管极大理想处的单个深度不等式成立。上个局部化命题要求 \(T_{\mathfrak m}\) 在它的全部支集素点满足条件，不能只检查其闭点。
\end{example}

零维环也说明这些条件不要求约化。例如 \(k[\varepsilon]/(\varepsilon^2)\) 只有一个素点，局部维数和深度均为零，所以满足所有 \(S_r\)。Serre 条件所检测的是局部支集维数与正则序列长度之间的关系；幂零结构、正则性和正规性还需要各自的条件。
% 纲要条目排为编号小节标题，并进入目录与书签；说明段仍保留于纲要注释中。

% !TeX root = ../../Book5.tex
% 纲要标题：### 8.4 极小自由消解和投射维数
% 纲要位置：几何学计划纲要.md，第 5327 行。
% 正文状态：初稿；尚未正式审核。

% 纲要内容开始
% 
% * 局部极小生成和极小微分
% * Betti 数及剩余域张量
% * 有限投射维数的检测
% * Auslander–Buchsbaum 公式的证明
% 
% 纲要内容结束

\section{极小自由消解和投射维数}
\label{g5:sec:0804}

自由消解的比较定理保证同调结果不依赖选择，但一个随意构造的消解可能含有许多可以消去的项。局部环上的 Nakayama 引理提供了逐层选取最少生成元的方法。这样得到的消解，在剩余域上所有微分都为零，各层自由秩因而成为模本身的不变量。本节最后把消解的长度与上一节的深度连接起来。

\subsection{局部极小生成和极小微分}
\label{g5:subsec:0804:01}

本节固定 Noether 局部环 \((A,\mathfrak m,k)\) 及非零有限 \(A\) 模 \(M\)。记
\(\mu_A(M)=\dim_k(M/\mathfrak mM)\)。由 Nakayama 引理，提升 \(M/\mathfrak mM\) 的一组基，便得到恰含 \(\mu_A(M)\) 个元素的生成系。它是极小生成系：少去任何一个元素，模 \(\mathfrak m\) 后便不能张成。

\begin{definition}\label{g5:def:minimal-free-resolution}
	设
	\[
		\cdots\longrightarrow F_2\xrightarrow{\mathrm d_2}F_1
		\xrightarrow{\mathrm d_1}F_0\xrightarrow{\epsilon}M\longrightarrow0
	\]
	是各项有限自由的消解。若对每个 \(i\geq1\) 都有 \(\mathrm d_i(F_i)\subseteq\mathfrak mF_{i-1}\)，则称它为\term[jixiaoziyouxiaojie]{极小自由消解}[minimal free resolution]。选定各项的基后，这等价于全部微分矩阵的元素都在 \(\mathfrak m\) 中。
\end{definition}

\begin{proposition}\label{g5:prop:minimal-resolution-existence}
	每个有限 \(A\) 模都有极小自由消解；其中 \(F_0\to M\) 及每个 \(F_i\to\ker d_{i-1}\) 都由目标的极小生成系给出，这里 \(\mathrm d_0=\epsilon\)。
\end{proposition}
\begin{proof}
	取 \(F_0=A^{\mu_A(M)}\)，用一组提升的基定义满射 \(\epsilon\)。它模 \(\mathfrak m\) 后是同构，故 \(K_1=\ker\epsilon\subseteq\mathfrak mF_0\)。Noether 性保证 \(K_1\) 有限，再取由极小生成系给出的满射 \(F_1\to K_1\)，与包含 \(K_1\hookrightarrow F_0\) 复合得到 \(\mathrm d_1\)。

	归纳地，若已取到满射 \(F_i\to K_i\)，令 \(K_{i+1}\) 为其核。模 \(\mathfrak m\) 后的映射 \(F_i/\mathfrak mF_i\to K_i/\mathfrak mK_i\) 为同构，所以 \(K_{i+1}\subseteq\mathfrak mF_i\)。取有限自由模 \(F_{i+1}\) 极小地满射到 \(K_{i+1}\)，继续即可。若某个核为零，后续所有项均取零。
\end{proof}

这种逐层构造还说明极小性的内容。任意各项有限自由的消解若满足定义中的条件，则 \(F_i\to K_i\) 的核包含于 \(\mathfrak mF_i\)，而张量的右正合性给出
\(F_i/\mathfrak mF_i\cong K_i/\mathfrak mK_i\)。所以每一层都恰好使用了所需的生成元数。这里的“极小”要求每一层的关系也极小，不能只检查 \(F_0\)。

\subsection{Betti 数及剩余域张量}
\label{g5:subsec:0804:02}

将极小自由消解与 \(k=A/\mathfrak m\) 张量，所有 \(\mathrm d_i\) 都变成零。消解原来记录的关系在剩余域中全部成为同调类，于是
\[
	\operatorname{Tor}_i^A(M,k)\cong F_i/\mathfrak mF_i.
\]
这既给出计算方法，也说明自由项的秩与选择无关。

\begin{definition}\label{g5:def:betti-number}
	模 \(M\) 的第 \(i\) 个 \term[Betti-shu]{Betti 数}[Betti number]\footnote{Enrico Betti（1823—1892），意大利数学家，研究代数方程、拓扑与弹性理论。}定义为
	\[
		\beta_i^A(M)=\dim_k\operatorname{Tor}_i^A(M,k).
	\]
	对极小自由消解，它等于 \(\operatorname{rank}_A F_i\)；特别地，\(\beta_0^A(M)=\mu_A(M)\)。
\end{definition}
\symidx{\(\beta_i^A(M)\)：局部有限模的第 \(i\) 个 Betti 数}

\begin{proposition}\label{g5:prop:minimal-resolution-uniqueness}
	同一个有限模的两个极小自由消解同构为带增广的复形。这个同构一般不是唯一的，但各层自由秩唯一。
\end{proposition}
\begin{proof}
	设两消解为 \(F_\bullet,G_\bullet\)。比较定理给出提升 \(\mathrm{id}_M\) 的链映射 \(f:F_\bullet\to G_\bullet\)、\(g:G_\bullet\to F_\bullet\)，并给出 \(gf\) 与 \(\mathrm{id}_F\)、\(fg\) 与 \(\mathrm{id}_G\) 之间的链同伦。模去 \(\mathfrak m\) 后，微分为零，同伦等式因而变成
	\[
		\overline g_i\overline f_i=\mathrm{id}_{F_i/\mathfrak mF_i},
		\qquad
		\overline f_i\overline g_i=\mathrm{id}_{G_i/\mathfrak mG_i}.
	\]
	故两边秩相同，\(f_i\) 是同阶方阵，且其行列式模 \(\mathfrak m\) 非零。在局部环中这个行列式是单位，所以 \(f_i\) 可逆。各 \(f_i\) 组成所需的复形同构。
\end{proof}

\begin{example}\label{g5:exm:dvr-betti}
	设 \(A\) 是 DVR，极大理想由 \(\pi\) 生成，\(r\geq1\)。模 \(A/(\pi^r)\) 的极小自由消解为
	\[
		0\longrightarrow A\xrightarrow{\ \pi^r\ }A
		\longrightarrow A/(\pi^r)\longrightarrow0.
	\]
	因此 \(\beta_0=\beta_1=1\)，其余 Betti 数为零；而其长度为 \(r\)。Betti 数记录各层最少生成元数，并不记录微分中 \(\pi\) 的幂次，也不能单独恢复模的长度。
\end{example}

\subsection{有限投射维数的检测}
\label{g5:subsec:0804:03}

\begin{definition}\label{g5:def:projective-dimension}
	非零 \(A\) 模 \(M\) 的\term[tousheweishu]{投射维数}[projective dimension]记为 \(\operatorname{pd}_A M\)，是使 \(M\) 存在长度为 \(n\) 的投射消解
	\[
		0\longrightarrow P_n\longrightarrow\cdots
		\longrightarrow P_0\longrightarrow M\longrightarrow0
	\]
	的最小整数 \(n\geq0\)。若不存在有限长度的投射消解，令其为 \(\infty\)。以下涉及深度公式时始终假设 \(M\ne0\)，不对零模代入数值等式。
\end{definition}
\symidx{\(\operatorname{pd}_A M\)：模的投射维数}

\begin{theorem}\label{g5:thm:projective-dimension-detection}
	设 \(M\) 为 Noether 局部环 \(A\) 上的非零有限模，\(n\geq0\)。以下条件等价：
	\begin{enumerate}
		\item \(\operatorname{pd}_A M\leq n\)；
		\item \(\operatorname{Tor}_{n+1}^A(M,k)=0\)；
		\item \(M\) 的极小自由消解满足 \(F_{n+1}=0\)；
		\item 对每个 \(A\) 模 \(N\)，均有 \(\operatorname{Ext}_A^{n+1}(M,N)=0\)。
	\end{enumerate}
	因此有限投射维数恰为极小自由消解最后一个非零项的次数，也等于最后一个非零 Betti 数的下标。
\end{theorem}
\begin{proof}
	第一条件给出长度不超过 \(n\) 的投射消解。由第一节的比较定理，投射消解同样可用于计算 Ext、Tor，故第二、第四条件成立。第二与第三条件等价，因为
	\(\operatorname{Tor}_{n+1}^A(M,k)\cong F_{n+1}/\mathfrak mF_{n+1}\)，而 \(F_{n+1}\) 有限自由。

	第三条件使极小消解在 \(F_n\) 处截断，因而推出第一条件。也可以直接在构造中看见这一点：\(F_{n+1}=0\) 表明第 \(n+1\) 个核为零，所有更高核及自由项都为零。

	最后，第四条件借助维数平移，等价于第 \(n\) 个合冲模 \(K_n\) 满足 \(\operatorname{Ext}_A^1(K_n,N)=0\) 对所有 \(N\) 成立；\(n=0\) 时取 \(K_0=M\)。由\cref{g5:prop:projectivity-ext}，\(K_n\) 投射。它又是有限模，所以由\cref{g5:prop:projective-local-free}知其自由。将前 \(n\) 步消解在 \(K_n\) 处截断，得到第一条件。
\end{proof}

这里一次正次数 Tor 的消失便迫使全部更高次数消失，关键在于底环局部、模有限，并且第二变量是剩余域。若随意换一个第二变量，便不能从这个证明得出同样的判断。有限性还保证 Nakayama 引理可以检测某一层确实为零。

\begin{example}\label{g5:exm:dual-number-infinite-pd}
	取 \(A=k[\varepsilon]/(\varepsilon^2)\)、\(M=k\)。第一节给出的每个微分均为乘 \(\varepsilon\) 的无限自由消解是极小的。于是
	\[
		\beta_i^A(k)=1\quad(i\geq0),\qquad
		\operatorname{pd}_A k=\infty.
	\]
	尽管 \(A\) 和 \(k\) 都是有限长度模，消解仍不会终止。这里 \(\varepsilon\) 的零化理想为极大理想，所以 \(\operatorname{depth}A=0\)；\(k\) 也有深度零。下一条公式中的有限投射维数假设因此不能省略。
\end{example}

\subsection{Auslander–Buchsbaum 公式的证明}
\label{g5:subsec:0804:04}

\term[Auslander-Buchsbaum-gongshi]{Auslander–Buchsbaum 公式}[Auslander--Buchsbaum formula]\footnote{Maurice Auslander（1926—1994），美国数学家，研究同调代数与表示论。David Alvin Buchsbaum（1929—2021），美国数学家，研究交换代数中的自由消解与同调维数。}把有限消解的长度解释为相对于底环损失的深度。证明的关键发生在长度一：极小消解中的矩阵元素属于 \(\mathfrak m\)，而 \(\mathfrak m\) 杀死 \(\operatorname{Ext}_A^i(k,A)\)，这会使一个关键映射等于零。

\begin{theorem}[Auslander–Buchsbaum 公式]\label{g5:thm:auslander-buchsbaum}
	设 \((A,\mathfrak m,k)\) 为 Noether 局部环，\(M\ne0\) 为有限 \(A\) 模。若 \(\operatorname{pd}_A M<\infty\)，则
	\[
		\operatorname{pd}_A M+\operatorname{depth}_A M
		=\operatorname{depth}A.
	\]
\end{theorem}
\begin{proof}
	置 \(d=\operatorname{depth}A\)、\(p=\operatorname{pd}_A M\)，对 \(p\) 归纳。若 \(p=0\)，则 \(M\) 有限投射，从而为正秩自由模。有限直和与 \(\operatorname{Ext}_A^i(k,-)\) 相容，所以 \(M\) 与 \(A\) 的首个非零 Ext 次数相同；由\cref{g5:thm:depth-ext}，\(\operatorname{depth}_A M=d\)。

	设 \(p=1\)。极小自由消解给出
	\[
		0\longrightarrow K\xrightarrow{u}F\longrightarrow M\longrightarrow0,
		\qquad K\cong A^a,\ F\cong A^b,\ a>0,
	\]
	且 \(u\) 的矩阵元素均在 \(\mathfrak m\) 中。由\cref{g5:prop:ext-annihilator}，\(\mathfrak m\operatorname{Ext}_A^i(k,A)=0\)；Ext 对第二变量的有限直和及标量映射的自然性表明
	\[
		\operatorname{Ext}_A^i(k,K)\xrightarrow{u_*}
		\operatorname{Ext}_A^i(k,F)
	\]
	在每个次数都为零。若 \(d=0\)，则 \(\operatorname{Hom}_A(k,K)\ne0\)，但它到 \(\operatorname{Hom}_A(k,F)\) 的映射一方面为零，另一方面因 Hom 左正合而单射，矛盾。因此 \(d\geq1\)。Ext 长正合列给出 \(j<d-1\) 时 \(\operatorname{Ext}_A^j(k,M)=0\)，并在首次可能非零的位置给出
	\[
		\operatorname{Ext}_A^{d-1}(k,M)
		\cong\operatorname{Ext}_A^d(k,K)\ne0.
	\]
	这里右端到 \(\operatorname{Ext}_A^d(k,F)\) 的映射为零，左端之前的 \(\operatorname{Ext}_A^{d-1}(k,F)\) 为零。于是 \(\operatorname{depth}_A M=d-1\)。

	现设 \(p\geq2\)。从极小消解取第一步
	\(0\to K\to F_0\to M\to0\)。它的尾部是 \(K\) 的极小消解，故 \(\operatorname{pd}_A K=p-1\)。归纳假设给出
	\[
		t:=\operatorname{depth}_A K=d-p+1<d.
	\]
	特别地 \(t\geq0\)，且 \(d\geq p-1\geq1\)。若 \(t=0\)，深度引理却给出
	\[
		0=\operatorname{depth}_A K
		\geq\min\{d,\operatorname{depth}_A M+1\}\geq1,
	\]
	矛盾。因此 \(t\geq1\)。对上述短正合列应用 \(\operatorname{Ext}_A^i(k,-)\)，利用 \(F_0\) 的深度为 \(d>t\)，得到
	\[
		\operatorname{Ext}_A^j(k,M)=0\quad(0\leq j<t-1),
		\qquad
		\operatorname{Ext}_A^{t-1}(k,M)
		\cong\operatorname{Ext}_A^t(k,K)\ne0.
	\]
	故 \(\operatorname{depth}_A M=t-1=d-p\)，归纳完成。
\end{proof}

公式说明，非零有限模一旦具有有限投射维数，该维数就不超过 \(\operatorname{depth}A\)。例如深度零的 Noether 局部环上，具有有限投射维数的非零有限模必定自由；上面的双数环却有投射维数无穷的剩余域。关于何种局部环能使全部有限模都具有有限投射维数，将在下一章结合正则性继续讨论。

\ProseParagraphBreak

本节的极小消解及公式证明可对照\cite[第19节，第153—155页]{Matsumura1989}。证明只用了本章已经建立的 Ext 长正合列、深度检测和极小性，没有预先调用正则局部环的同调刻画。
% 纲要条目排为编号小节标题，并进入目录与书签；说明段仍保留于纲要注释中。

% !TeX root = ../../Book5.tex
% 纲要标题：### 8.5 交与嵌入的模型
% 纲要位置：几何学计划纲要.md，第 5334 行。
% 正文状态：初稿；尚未正式审核。

% 纲要内容开始
% 
% * 横截和非横截的 Tor 对照
% * 超曲面与完全交
% * 嵌入分支导致的深度下降
% * 同一支集的不同局部同调信息
% 
% ---
% 
% 纲要内容结束

\section{交与嵌入的模型}
\label{g5:sec:0805}

前面的理论给出两类可直接计算的资料：把消解与另一个模张量，得到交中的 Tor；把极小消解与剩余域张量，得到 Betti 数与投射维数。本节在平面上的几个局部模型中实际求出这些资料，同时比较支集、交的长度和深度。固定一个域 \(k\)，主要使用
\(R=k[x,y]_{(x,y)}\)，其极大理想为 \(\mathfrak n=(x,y)R\)。

\subsection{横截和非横截的 Tor 对照}
\label{g5:subsec:0805:01}

在同一仿射空间中由理想 \(I,J\) 给出的两个闭子对象，其共同方程为 \(I+J\)。代数上有
\[
	(R/I)\otimes_R(R/J)\cong R/(I+J).
\]
第一节又从短正合列 \(0\to I\to R\to R/I\to0\) 得到
\[
	\operatorname{Tor}_1^R(R/I,R/J)
	\cong\frac{I\cap J}{IJ}.
\]
右边比较“同时满足两个理想条件”与“由两个理想的乘积产生”的差异。更高次的 Tor 则需要继续追踪关系之间的关系。

\begin{example}\label{g5:exm:transverse-axis-tor}
	取 \(M=R/(x)\)、\(N=R/(y)\)。前者的自由消解为
	\[
		0\longrightarrow R\xrightarrow{x}R\longrightarrow M\longrightarrow0.
	\]
	与 \(N\) 张量后，正次数部分只有 \(N\xrightarrow{x}N\)。由于
	\(N\cong k[x]_{(x)}\)，乘 \(x\) 单射，故
	\[
		\operatorname{Tor}_i^R(M,N)=0\quad(i>0),\qquad
		M\otimes_RN\cong k.
	\]
	两条坐标轴在原点的切方向不同，且交的局部代数长度为一。这是横截相交的基本模型。
\end{example}

\begin{example}\label{g5:exm:tangent-curves-tor}
	取 \(M=R/(y)\)、\(N=R/(y-x^r)\)，其中 \(r\geq2\)。两条曲线在原点的切线都是 \(y=0\)。用 \(M\) 的一步自由消解张量 \(N\)，所得微分为乘 \(y\)。在同构 \(N\cong k[x]_{(x)}\) 下，它就是乘 \(x^r\)，仍然单射。因此
	\[
		\operatorname{Tor}_i^R(M,N)=0\quad(i>0),\qquad
		M\otimes_RN\cong k[x]_{(x)}/(x^r).
	\]
	这次交的局部代数长度为 \(r\)。相切已经反映在零次张量积的幂零结构和长度中，正次数 Tor 却全部消失。即使两条曲线本身光滑，Tor 的消失也不单独蕴含切空间意义的横截。
\end{example}

\begin{example}\label{g5:exm:self-intersection-tor}
	取 \(M=N=R/(x)\)，即同一条直线与自身相交。把乘 \(x\) 的消解与 \(M\) 张量，微分变为零，故
	\[
		\operatorname{Tor}_0^R(M,M)\cong M,\qquad
		\operatorname{Tor}_1^R(M,M)\cong M,\qquad
		\operatorname{Tor}_i^R(M,M)=0\quad(i\geq2).
	\]
	这里共同方程只切掉一个方向，交仍是一条线。第一 Tor 保留了重复施加同一条方程所产生的关系；它也等于 \((x)/(x^2)\)。
\end{example}

三例分别给出长度一的孤立交、带高阶接触的孤立交和具有共同分支的交。把正次数 Tor 全部为零的条件称为两个模的\term[Tor-duli]{Tor 独立}[Tor-independence]。它描述张量消解的正合性；判断几何横截时还要使用切空间或微分条件。今后引入交数时，零次张量的长度与高次 Tor 的长度将各有作用。

\subsection{超曲面与完全交}
\label{g5:subsec:0805:02}

设 \((A,\mathfrak m,k)\) 为 Noether 局部环。第三章已把正则序列给出的局部方程作为完全交模型。现在 Koszul 消解能同时给出它的 Betti 数、投射维数和相对于环境环的深度。

\begin{proposition}\label{g5:prop:complete-intersection-homology}
	设 \(f_1,\ldots,f_c\in\mathfrak m\) 为 \(A\) 正则序列，\(I=(f_1,\ldots,f_c)\)、\(B=A/I\)。则
	\[
		\beta_i^A(B)=\binom ci\quad(0\leq i\leq c),\qquad
		\operatorname{pd}_A B=c,\qquad
		\operatorname{depth}_A B=\operatorname{depth}A-c.
	\]
	此外，作为 \(B\) 模，
	\[
		\operatorname{Tor}_i^A(B,B)
		\cong\bigwedge\nolimits_B^i B^c\quad(0\leq i\leq c),
	\]
	且 \(i>c\) 时为零。这里 \(B^c\) 的基来自所选方程的次序。
\end{proposition}
\begin{proof}
	由\cref{g5:thm:regular-koszul-acyclic}，\(K_\bullet(f_1,\ldots,f_c;A)\) 是 \(B\) 的自由消解，第 \(i\) 项秩为 \(\binom ci\)。微分矩阵的元素是各 \(f_j\) 及其相反数，都属于 \(\mathfrak m\)，故它极小。其顶次项为 \(A\ne0\)，所以投射维数等于 \(c\)。深度等式由 Auslander–Buchsbaum 公式得到，也可逐次应用正则元的深度下降公式。

	再与 \(B\) 张量，各 \(f_j\) 的像都为零，因此全部微分为零；第 \(i\) 项本身就是同调，得到最后的同构。
\end{proof}

当 \(c=1\) 时，非零因子 \(f\) 给出一个超曲面局部模型 \(A/(f)\)。它相对于环境环只有两项自由消解，深度下降一。若 \(f=x^2\in R\)，得到双直线 \(R/(x^2)\)：它含有非零幂零元 \(\overline x\)，但 \(y\) 仍为非零因子，深度恰为一。若取正则序列 \(x,y\)，则
\[
	0\longrightarrow R\xrightarrow{\binom{-y}{x}}R^2
	\xrightarrow{(x\ \ y)}R\longrightarrow k\longrightarrow0
\]
是原点的极小自由消解，Betti 数为 \(1,2,1\)。

\ProseParagraphBreak

上述投射维数的下标 \(A\) 是必需的。例如 \(B=R/(x^2)\) 作为 \(R\) 模的投射维数为一，但作为 \(B\) 模它自由，投射维数为零；而 \(B/(x)\) 作为 \(B\) 模有各微分均为乘 \(\overline x\) 的无限极小自由消解。因为在 \(B\) 中 \(\operatorname{Ann}_B(\overline x)=(\overline x)\)，该无限复形确实正合，且每层都非零。环境环改变时，所要测量的关系也改变。

\subsection{嵌入分支导致的深度下降}
\label{g5:subsec:0805:03}

第二章的嵌入点模型现在可以逐层计算。取
\[
	E=R/(x^2,xy).
\]
其支集仍是直线 \(V(x)\)，但非零元素 \(\overline x\) 同时被 \(x,y\) 杀死。因此 \(\operatorname{Hom}_R(k,E)\ne0\)，深度检测立即给出 \(\operatorname{depth}_R E=0\)。这个零维子模位于原点，正对应于嵌入的关联素理想 \(\mathfrak n\)。

\begin{proposition}\label{g5:prop:embedded-point-resolution}
	模 \(E\) 有极小自由消解
	\[
		0\longrightarrow R
		\xrightarrow{\binom{-y}{x}}R^2
		\xrightarrow{(x^2\ \ xy)}R
		\longrightarrow E\longrightarrow0.
	\]
	故其 Betti 数为 \(1,2,1\)，投射维数为二，深度为零。
\end{proposition}
\begin{proof}
	两个矩阵的乘积为零。若 \((a,b)\) 属于中间微分的核，则
	\(x(xa+yb)=0\)。环 \(R\) 为整环，故 \(xa+yb=0\)。模去 \(x\) 后，\(yb=0\) 在 \(R/(x)\cong k[y]_{(y)}\) 中成立，因此 \(b=xc\)。代回得 \(a=-yc\)，所以核恰为列向量 \((-y,x)^{\mathsf T}\) 的像。左端映射单射，因为 \(xc=0\) 推出 \(c=0\)。最右端的像为定义理想，故整列正合。

	所有矩阵元素都在 \(\mathfrak n\) 中，消解极小，因此 Betti 数及投射维数如述。环 \(R\) 的维数为二，且 \(x,y\) 正则，所以 \(\operatorname{depth}R=2\)；Auslander–Buchsbaum 公式也给出 \(\operatorname{depth}_R E=2-2=0\)。
\end{proof}

在泛点 \((x)\) 处局部化，\(y\) 成为单位，关系 \(xy=0\) 使 \(x=0\)，于是 \(E_{(x)}\cong(R/(x))_{(x)}\)。在闭点，子模 \(R\overline x\cong k\) 却保留下来；具体地有短正合列
\[
	0\longrightarrow k\xrightarrow{1\mapsto\overline x}
	E\longrightarrow R/(x)\longrightarrow0.
\]
其中 \(\overline x\ne0\)，因为 \(x\in(x^2,xy)=x\mathfrak n\) 会在整环中推出 \(1\in\mathfrak n\)。这一列同时显示出直线部分和嵌入点部分。

\ProseParagraphBreak

与双直线 \(R/(x^2)\) 比较，二者都有非零幂零元和一维支集，但前者的 \(y\) 是非零因子，后者的 \(y\) 杀死 \(\overline x\)。使深度降到零的具体原因，是出现被整个极大理想杀死的非零元素。不能只凭“有幂零元”便断定深度下降到零。

\subsection{同一支集的不同局部同调信息}
\label{g5:subsec:0805:04}

支集只回答模在哪些点非零。即使支集和零化理想都相同，深度与 Betti 数仍可能不同。取
\[
	P=R,\qquad Q=R\oplus k.
\]
二者的零化理想均为零，支集均为整个 \(\operatorname{Spec}R\)，在泛点的秩也都为一。可是 \(Q\) 中的第二个直和项使 \(\operatorname{Hom}_R(k,Q)\ne0\)，故
\[
	\operatorname{depth}_R P=2,\qquad
	\operatorname{depth}_R Q=0.
\]

\ProseParagraphBreak

把 \(k\) 的 Koszul 消解与只在零次有一项 \(R\) 的复形直和，得到 \(Q\) 的极小自由消解
\[
	0\longrightarrow R
	\xrightarrow{\binom{-y}{x}}R^2
	\xrightarrow{\left(\begin{smallmatrix}0&0\\ x&y\end{smallmatrix}\right)}
	R^2\longrightarrow Q\longrightarrow0.
\]
增广为 \((a,b)\mapsto(a,\overline b)\)。因此
\[
	(\beta_0^R(P),\beta_1^R(P),\beta_2^R(P))=(1,0,0),
	\qquad
	(\beta_0^R(Q),\beta_1^R(Q),\beta_2^R(Q))=(2,2,1).
\]
后一模的所有额外关系都集中在闭点：若 \(\mathfrak p\ne\mathfrak n\)，可取 \(a\in\mathfrak n\setminus\mathfrak p\)；它在 \(R_{\mathfrak p}\) 中可逆却杀死 \(k\)，故 \(k_{\mathfrak p}=0\) 且 \(Q_{\mathfrak p}\cong R_{\mathfrak p}\)。

\ProseParagraphBreak

在闭点，\(Q\) 的局部维数为二而深度为零，因此它不满足 \(S_1\)。这个例子把三种资料分开：支集描述非零位置，泛秩描述稠密处的自由部分，局部同调则能发现某一点新增的生成元、关系与嵌入子模。下一章将用深度和消解讨论局部环的正则性及深度等于维数的情形，把这些计算组织成结构判据。
% 纲要条目排为编号小节标题，并进入目录与书签；说明段仍保留于纲要注释中。


\begin{chaptersummary}
	自由消解通过逐层记录关系把模变成复形。比较定理和链同伦使 Ext、Tor 与选择无关，长正合列则将一次正合性的缺失转移到相邻次数。合冲模和维数平移使高次问题可以回到一次提升或一次张量单射。

	\ProseParagraphBreak

	Koszul 复形由外代数及乘方程的微分构成。正则序列保证正次数同调消失；对 Noether 局部环的非零有限模和极大理想中的序列，第一同调消失已经足以反推正则性。末端商非零、局部性和有限性都有实际作用，换序与基变换须在这些条件下使用。

	\ProseParagraphBreak

	深度等于极大正则序列的共同长度，也等于 \(\operatorname{Ext}_A^i(k,M)\) 首次非零的次数，并且不超过模的局部维数。极小自由消解的微分在剩余域上归零，其自由秩便成为 Betti 数。有限投射维数时，消解长度与模的深度之和等于底环深度。具体的交与嵌入点计算表明，支集、交的长度和局部同调提供相互补充的资料。
\end{chaptersummary}

\BookExercises

\begin{exercise}\label{g5:ex:integer-tor-ext}
	设 \(m,n\geq2\)，\(g=\gcd(m,n)\)。计算
	\(\operatorname{Tor}_i^{\mathbb Z}(\mathbb Z/n\mathbb Z,\mathbb Z/m\mathbb Z)\) 和
	\(\operatorname{Ext}_{\mathbb Z}^i(\mathbb Z/n\mathbb Z,\mathbb Z/m\mathbb Z)\) 的全部次数。
\end{exercise}
\begin{solution}
	用自由消解 \(0\to\mathbb Z\xrightarrow{n}\mathbb Z\to\mathbb Z/n\mathbb Z\to0\)。
	令 \(N=\mathbb Z/m\mathbb Z\)，则零次 Tor 和一次 Ext 都为 \(N/nN\)，一次 Tor 和零次 Ext 都为 \(N[n]=\{a\in N:na=0\}\)。其余正高次均为零。
	在 \(N\) 中，\(nN=gN\)，所以 \(N/nN\cong\mathbb Z/g\mathbb Z\)；而 \(na=0\) 等价于 \(m/g\) 整除 \(a\)，故 \(N[n]\) 由 \(m/g\) 的类生成，也同构于 \(\mathbb Z/g\mathbb Z\)。因此次数零、一的四个群都同构于 \(\mathbb Z/g\mathbb Z\)，次数至少二时全为零。
\end{solution}

\begin{exercise}\label{g5:ex:ext-extension-principal}
	设 \(f\) 为环 \(A\) 的非零因子，\(N\) 为任意 \(A\) 模。证明
	\(\operatorname{Ext}_A^1(A/(f),N)\cong N/fN\)，并把右边元素解释为短正合列
	\(0\to N\to E\to A/(f)\to0\) 的等价类；等价要求在两端诱导恒等映射。
\end{exercise}
\begin{solution}
	对一步自由消解应用 \(\operatorname{Hom}_A(-,N)\)，得到 \(N\xrightarrow{f}N\)，其一次上同调为 \(N/fN\)。

	在给定扩张中，选取 \(1\bmod f\) 的提升 \(e\in E\)。元素 \(fe\) 落在 \(N\) 中，记为 \(n\)。把 \(e\) 改成 \(e+t\) 会把 \(n\) 改成 \(n+ft\)，所以扩张确定一个 \(N/fN\) 中的类。反过来，对 \(n\in N\) 令
	\[
		E_n=(N\oplus A)/A(-n,f).
	\]
	映射 \(N\to E_n\)、\(u\mapsto[(u,0)]\) 单射：若 \((u,0)=a(-n,f)\)，则 \(af=0\)，由非零因子条件得 \(a=0\)。映射 \(E_n\to A/(f)\)、\([(u,a)]\mapsto a\bmod f\) 满射；若 \(a=fb\)，则 \([(u,a)]=[(u+bn,0)]\)，所以其核正是 \(N\)。取 \(e=[(0,1)]\)，有 \(fe=n\)。

	对任意给定扩张，\((u,a)\mapsto u+ae\) 诱导 \(E_n\to E\)。它满射；若 \(u+ae=0\)，则 \(a=fb\) 且 \(u=-bn\)，故核恰为 \(A(-n,f)\)，从而是同构。改变提升只改变 \(n\) 的模 \(fN\) 类，相反同类元素可通过这个改变构造扩张等价。因此上述对应双射，且零类恰对应可选择 \(fe=0\) 的分裂扩张。
\end{solution}

\begin{exercise}\label{g5:ex:koszul-embedded-homology}
	在 \(R=k[x,y]_{(x,y)}\) 中，计算 \(K_\bullet(x^2,xy;R)\) 的全部同调，并据此判断该序列是否正则。
\end{exercise}
\begin{solution}
	微分为
	\[
		0\longrightarrow R\xrightarrow{\binom{-xy}{x^2}}R^2
		\xrightarrow{(x^2\ \ xy)}R\longrightarrow0.
	\]
	零次同调是 \(R/(x^2,xy)\)。环为整环，所以最左端映射单射，二次同调为零。若 \(x^2a+xyb=0\)，则 \(xa+yb=0\)，与正文相同的模 \(x\) 消去给出 \((a,b)=c(-y,x)\)。边界恰为 \(xR(-y,x)\)，所以
	\(H_1\cong R/(x)\ne0\)。故序列不正则。直接检查也可看见：在 \(R/(x^2)\) 中，非零元素 \(\overline x\) 被 \(xy\) 杀死。
\end{solution}

\begin{exercise}\label{g5:ex:koszul-unit-contraction}
	若序列 \(a_1,\ldots,a_r\) 中有一个单位，证明其 Koszul 复形在所有次数都同调为零，并说明为何不能据此称该序列正则。
\end{exercise}
\begin{solution}
	设 \(a_j\) 可逆。令 \(h(\omega)=a_j^{-1}e_j\wedge\omega\)。Koszul 微分的交错公式给出
	\(\mathrm dh+h\mathrm d=\mathrm{id}\)，所以每个循环 \(\omega\) 都等于 \(\mathrm d(h\omega)\)，全部同调为零。另一方面，\((a_1,\ldots,a_r)=A\)，末端商为零，违反正则序列定义中的末端非零条件。局部反向判据要求各 \(a_i\) 属于极大理想，正是为了排除这种情形。
\end{solution}

\begin{exercise}\label{g5:ex:module-regular-not-ring}
	取 \(A=(k[x,y]/(xy))_{(x,y)}\)、\(M=A/(y)\)。证明 \(x\) 是 \(M\) 正则元，却不是 \(A\) 正则元；分别计算 \(K_\bullet(x;M)\) 与 \(K_\bullet(x;A)\) 的一次同调。
\end{exercise}
\begin{solution}
	有 \(M\cong k[x]_{(x)}\)，乘 \(x\) 单射且 \(M/xM\cong k\ne0\)，所以 \(x\) 为 \(M\) 正则元，\(H_1(x;M)=0\)。
	在 \(A\) 中，\(y\ne0\) 且 \(xy=0\)，故 \(x\) 不是非零因子。其零化理想恰为 \((y)A\)：在 \(k[x,y]_{(x,y)}\) 中，若 \(xu\in(xy)\)，可约去 \(x\) 得 \(u\in(y)\)。因此 \(H_1(x;A)=(y)A\ne0\)。
\end{solution}

\begin{exercise}\label{g5:ex:koszul-base-change-limits}
	以 \(A=k[t]\) 和序列 \(t\) 为例，分别考察基变换 \(A\to A/(t)\) 与 \(A\to A_t\)。说明前者如何破坏正次数正合性，后者为何保持正合性却不保持正则序列的末端条件。
\end{exercise}
\begin{solution}
	原复形为 \(0\to A\xrightarrow{t}A\to0\)，其正次数同调为零。
	对 \(B=A/(t)=k\)，张量后的微分为零，故 \(H_1=k\ne0\)。这次基变换不平坦。
	对 \(B=A_t\)，乘 \(t\) 可逆，所以零次、一次同调都为零。局部化平坦，确实保持正次数正合性；但 \(B/tB=0\)，因此像序列不再满足末端商非零。平坦基变换还须保留末端非零模，才能推出像序列正则；忠实平坦性保证这一点。
\end{solution}

\begin{exercise}\label{g5:ex:monomial-complete-intersection}
	设 \(R=k[x,y]_{(x,y)}\)、\(a,b\geq1\)、\(L=R/(x^a,y^b)\)。证明 \(x^a,y^b\) 正则，求 \(L\) 的长度、Betti 数、投射维数和深度。
\end{exercise}
\begin{solution}
	乘 \(x^a\) 在整环 \(R\) 中单射。在 \(k[x,y]/(x^a)\) 中，各元素可唯一写成 \(\sum_{i=0}^{a-1}x^ip_i(y)\)，故乘 \(y^b\) 单射；局部化保持该单射。末端商非零，因而序列正则。
	商中的单项式 \(x^iy^j\)、\(0\leq i<a,\ 0\leq j<b\)，构成 \(k\) 基。这个商只有极大理想 \((x,y)\)，再局部化不会改变它，所以长度等于其 \(k\) 维数 \(ab\)。Koszul 消解极小，故 Betti 数为 \(1,2,1\)，投射维数二。环 \(R\) 深度二，Auslander–Buchsbaum 公式给出 \(\operatorname{depth}_R L=0\)，也可直接用被 \((x,y)\) 杀死的非零元素 \(x^{a-1}y^{b-1}\) 检测。
\end{solution}

\begin{exercise}\label{g5:ex:depth-direct-sum}
	证明有限个非零有限模的直和，其深度为各项深度的最小值。进而对 \(R=k[x,y]_{(x,y)}\)、
	\(M=R/(x)\oplus R/(y)\)，求深度与维数，并判断 \(S_1\) 条件。
\end{exercise}
\begin{solution}
	固定剩余域 \(k\) 的自由消解，对第二变量取有限直和时，Hom 复形逐项分解，其上同调也分解。因此 Ext 首次非零的次数是各直和项对应次数的最小值，深度结论由 Ext 刻画得到。
	在 \(R/(x)\) 上 \(y\) 正则，在 \(R/(y)\) 上 \(x\) 正则，二者维数都为一，故深度都为一。因此 \(M\) 深度一，支集为两条轴的并，维数也为一。支集中的素理想只有 \((x)\)、\((y)\)、\((x,y)\)：前两个点处局部化的模维数为零，\(S_1\) 不再要求正深度；闭点处深度与维数均为一。因此 \(M\) 满足 \(S_1\)。
\end{solution}

\begin{exercise}\label{g5:ex:maximal-ideal-depth}
	设 \((A,\mathfrak m,k)\) 为深度正的 Noether 局部环。证明 \(\operatorname{depth}_A\mathfrak m=1\)。取 \(A=k[x,y,z]_{(x,y,z)}\) 时，求 \(\mathfrak m\) 的 Betti 数及投射维数。
\end{exercise}
\begin{solution}
	深度正使 \(\operatorname{Hom}_A(k,A)=0\)，且 \(A\) 不是域，故 \(\mathfrak m\ne0\)。从 \(0\to\mathfrak m\to A\to k\to0\) 的 Ext 长正合列得到
	\[
		\operatorname{Hom}_A(k,\mathfrak m)=0,\qquad
		0\longrightarrow\operatorname{Hom}_A(k,k)
		\longrightarrow\operatorname{Ext}_A^1(k,\mathfrak m).
	\]
	右边非零，所以深度恰为一。
	在三变量例子中，\(x,y,z\) 是正则序列；将 \(k\) 的 Koszul 消解去掉零次的 \(A\to k\)，其像为 \(\mathfrak m\)，得到极小消解
	\[
		0\longrightarrow A\longrightarrow A^3
		\longrightarrow A^3\longrightarrow\mathfrak m\longrightarrow0.
	\]
	所以 Betti 数为 \(3,3,1\)，投射维数为二，与 \(2+1=\operatorname{depth}A=3\) 相符。
\end{solution}

\begin{exercise}\label{g5:ex:completion-betti}
	设 \((A,\mathfrak m,k)\) 为 Noether 局部环，\(M\ne0\) 为有限模。证明
	\[
		\beta_i^A(M)=\beta_i^{\widehat A}(\widehat M)
		\quad(i\geq0),\qquad
		\operatorname{pd}_A M=\operatorname{pd}_{\widehat A}\widehat M,
	\]
	后一个等式允许两边同时无穷。
\end{exercise}
\begin{solution}
	第七章已证明 \(\widehat A\) 为 Noether 局部环，极大理想为 \(\mathfrak m\widehat A\)，剩余域为 \(k\)，且对有限模有 \(\widehat M\cong M\otimes_A\widehat A\)。取 \(M\) 的极小自由消解，与平坦的 \(\widehat A\) 张量，仍然正合；每个自由项的秩不变，微分矩阵元素落在 \(\mathfrak m\widehat A\) 中，故所得消解极小。忠实平坦性保证 \(\widehat M\ne0\)。因此各 Betti 数相等，最后非零项的次数也相等；若原消解永不终止，新消解同样不会终止。
\end{solution}

\begin{exercise}\label{g5:ex:quotient-ring-depth}
	设 \((A,\mathfrak m)\) 为 Noether 局部环，\(I\subseteq\mathfrak m\)，\(B=A/I\)，\(M\ne0\) 为有限 \(B\) 模。证明 \(\operatorname{depth}_A M=\operatorname{depth}_B M\)，并用 \(M=B=A/(f)\) 说明投射维数一般不满足相同的等式，其中 \(f\in\mathfrak m\) 是非零因子。
\end{exercise}
\begin{solution}
	每个 \(a\in\mathfrak m\) 在 \(M\) 上的作用只依赖其像 \(\overline a\in\mathfrak m/I\)。逐次商去前面元素后仍然如此，所以一个 \(\mathfrak m\) 中的序列是 \(M\) 正则序列，当且仅当其像序列在 \(B\) 上正则；非零的末端商在两个环下是同一个模。反过来，\(B\) 中的每个序列都可逐项提升到 \(\mathfrak m\)。因此可取得的长度集合相同，深度相等。
	若 \(M=B=A/(f)\)，其关于 \(A\) 的一步自由消解极小，故 \(\operatorname{pd}_A B=1\)；关于 \(B\)，它是秩一自由模，故 \(\operatorname{pd}_B B=0\)。深度在这个比较中相同，底环的深度却改变了，两个投射维数分别适用于各自底环。
\end{solution}

\begin{exercise}\label{g5:ex:squared-maximal-ideal-resolution}
	在 \(R=k[x,y]_{(x,y)}\) 中，为 \(R/(x,y)^2\) 构造极小自由消解，并求其长度、Betti 数与深度。
\end{exercise}
\begin{solution}
	所需消解为
	\[
		0\longrightarrow R^2
		\xrightarrow{\left(\begin{smallmatrix}-y&0\\x&-y\\0&x\end{smallmatrix}\right)}
		R^3\xrightarrow{(x^2\ \ xy\ \ y^2)}R
		\longrightarrow R/(x,y)^2\longrightarrow0.
	\]
	验证中间的核：若 \(x^2a+xyb+y^2c=0\)，模去 \(x\) 得 \(c=xc_1\)；约去 \(x\) 后得 \(xa+yb+y^2c_1=0\)。再模去 \(y\) 得 \(a=ya_1\)，代回得到 \(b=-xa_1-yc_1\)。所以
	\[
		(a,b,c)=-a_1(-y,x,0)+c_1(0,-y,x).
	\]
	这两个列向量之间没有非平凡关系：第一坐标先使第一个系数为零，第三坐标再使第二个系数为零，因为 \(R\) 是整环。故整列正合，且所有矩阵元素都在极大理想中，消解极小。Betti 数为 \(1,3,2\)，投射维数二。商的 \(k\) 基为 \(1,\overline x,\overline y\)，所以长度三；\(\overline x\) 被整个极大理想杀死，深度为零。
\end{solution}

\begin{exercise}\label{g5:ex:alternating-betti-rank}
	设 \(A\) 为 Noether 局部整环，\(M\ne0\) 为有限模且 \(\operatorname{pd}_A M=p<\infty\)。证明
	\[
		\operatorname{rank}_A M=\sum_{i=0}^p(-1)^i\beta_i^A(M).
	\]
	用嵌入点模 \(R/(x^2,xy)\) 和 \(R\oplus k\) 核对该等式。
\end{exercise}
\begin{solution}
	与分式域 \(K\) 张量保持极小自由消解的正合性，得到有限维向量空间的有限正合列。逐次用短正合列的维数加性，得
	\(\dim_K(M\otimes_AK)=\sum_i(-1)^i\dim_K(F_i\otimes_AK)\)，即所求等式。
	嵌入点模被非零元素 \(x^2\) 杀死，故秩零，其 Betti 数满足 \(1-2+1=0\)。
	模 \(R\oplus k\) 的泛秩为一，其 Betti 数满足 \(2-2+1=1\)。交错和给出泛秩，但不保留各层的全部关系数。
\end{solution}

\begin{exercise}\label{g5:ex:double-line-residue-resolution}
	令 \(B=(k[x,y]/(x^2))_{(x,y)}\)，仍用 \(x,y\) 表示它们的像。证明其剩余域 \(k\) 有极小自由消解
	\[
		\cdots\xrightarrow{D_-}B^2\xrightarrow{D_+}B^2
		\xrightarrow{D_-}B^2\xrightarrow{(x\ \ y)}B
		\longrightarrow k\longrightarrow0,
		\qquad
		D_\pm=\begin{pmatrix}x&\pm y\\0&x\end{pmatrix}.
	\]
	求全部 Betti 数及 \(\operatorname{pd}_B k\)。
\end{exercise}
\begin{solution}
	在 \(B\) 中，\(\operatorname{Ann}(x)=(x)\)，且 \(B/(x)\cong k[y]_{(y)}\)，乘 \(y\) 单射。若 \(xa+yb=0\)，模去 \(x\) 得 \(b=xc\)，于是 \(x(a+yc)=0\)，可写 \(a=-yc+xd\)。故
	\((a,b)=D_-(d,c)\)，给出零次之前的正合性。

	若 \(D_-(a,b)=0\)，则 \(xb=0\)，可写 \(b=xc\)；第一行给出 \(x(a-yc)=0\)，故 \(a=yc+xd\)。因此 \((a,b)=D_+(d,c)\)。反过来 \(D_-D_+=0\)，所以 \(\ker D_-=\operatorname{im}D_+\)。
	类似地，\(D_+(a,b)=0\) 使 \(b=xc\)、\(a=-yc+xd\)，故 \(\ker D_+=\operatorname{im}D_-\)。所有次数都正合；这个验证也适用于特征二，因为只用了加法逆元和 \(x^2=0\)。

	矩阵元素均在极大理想中，所以消解极小。于是
	\(\beta_0^B(k)=1\)、\(\beta_i^B(k)=2\) 对所有 \(i\geq1\) 成立，投射维数为无穷。另一方面，\(y\) 是 \(B\) 的非零因子且 \(\dim B=1\)，故 \(\operatorname{depth}B=1\)，而 \(\operatorname{depth}_B k=0\)。深度差为一并不能单独推出投射维数有限。
\end{solution}
